\documentclass[11pt]{article}
\usepackage[hmargin=.9in,vmargin=.82in]{geometry}
\usepackage{amsmath,amssymb,amsthm,booktabs,mathtools,needspace}
\usepackage{pgfplots}
\pgfplotsset{compat=1.18}
\usepackage[hidelinks]{hyperref}
\hypersetup{
  pdftitle={Finite-Time Observability in Gowdy Cosmology},
  pdfauthor={Ronald Bibb},
  pdfsubject={Exact collisions, stable recovery, and the metric modulus}
}
\newtheorem{theorem}{Theorem}
\newtheorem{proposition}{Proposition}
\newcommand{\sech}{\operatorname{sech}}
\newcommand{\TT}{\mathbb{T}}
\makeatletter
\def\input@path{{./}{../}}
\makeatother
\title{Finite-Time Observability in Gowdy Cosmology\protect\\
Exact Collisions, Stable Recovery, and the Metric Modulus}
\author{Ronald Bibb}
\date{Working draft, 25 November 2026}
\begin{document}
\maketitle

\begin{abstract}
We study recovery of missing finite-time velocity from Gowdy
snapshots. For homogeneous vacuum solutions with speed at most
$V$, the mixed record is injective exactly when
$V\log(b/a)\le 2.399357\ldots$. Exact inhomogeneous collisions
can match both induced metrics and the known velocity component
at resonant slices, even when both target directions are active.
For these rank-two solutions a suitably placed full orbit
snapshot instead gives a local Lipschitz inverse. Relative
parallel transport along the two null directions explains
the gain and gives stable fixed-slice recovery near
inhomogeneous backgrounds. Closed ray transports give
a separate criterion when neither velocity is supplied.
The metric gradient can restore an energy estimate at
irrational time gaps where orbit fields alone fail.
In an open low-velocity regime, the regular-time inverse
determines a Lipschitz asymptotic state.
Identical measured metrics can also hide differences in
the leading speed or, at the same speed, in the limiting
metric potential. A frequency-assigned
alternative is stable uniformly in bandwidth on short windows
of nonlinear evolution. For the archived spike-region restart
we certify a finite-grid nonlinear contraction from the original
secant. For polarized Einstein--scalar data we derive a sharp
logarithmic regularity threshold for estimating the limiting
metric modulus from two field snapshots and one metric anchor.
\end{abstract}

\section{The observation map}
A finite-window logarithmic slope and intercept of $P$, together
with endpoint $Q$ and $e^{2P}Q_\tau$ in log time $\tau=-\log t$,
are algebraically equivalent to the record
\[
 \mathcal O_0=(P(a),Q(a),Q_t(a),P(b)),\qquad 0<a<b.
\]
We may also retain the additive metric anchor $\lambda(a,0)$.
The original finite-window study measured a large error when
$P_\tau(a)$ was replaced with a secant \cite{old}.
A later computation recovered the tested spike-region history with
the same data by solving the PDE over the observation window.
This raises a mathematical question about the observation map itself.
Two-time observation and small-denominator results for scalar waves
give an important comparison \cite{vest,cssw}. Gowdy forward and
backward computation is established \cite{amorim}. We examine what
changes when the two field channels of the vacuum system interact.
The sampling construction below is related to dynamical
sampling and discrete-time wave data assimilation
\cite{aldroubi,cindea}. We use full spatial Fourier coefficients
at a choice of times, with a nonlinear Cauchy state hidden
only in its initial velocity. The null-transport argument
is related to geometric-optics observability for coupled
waves \cite{dlr,clw}, here with discrete orbital snapshots
rather than observation through a spatial control window.

Write $U=(P,Q)$ as a wave map to the upper half-plane with metric
$dP^2+e^{2P}dQ^2$. The equations in areal time are
\begin{equation}
 D_t U_t+t^{-1}U_t-D_\theta U_\theta=0.
 \label{wave}
\end{equation}
We take $\theta\in\TT=\mathbb R/(2\pi\mathbb Z)$ and use a fixed
marking and Killing basis for the measured fields.
The remaining metric potential is determined up to an additive
constant by constraints, with
$\lambda_t=-t(|U_t|^2+|U_\theta|^2)$ and
$\lambda_\theta=-2t\langle U_t,U_\theta\rangle$ in the
positive-$\lambda_\tau$ convention \cite{amorim}.
In this convention the four-metric is
\[
 ds^2=t^{-1/2}e^{-\lambda/2}(-dt^2+d\theta^2)
       +t\big[e^P(dx+Q\,dy)^2+e^{-P}dy^2\big].
\]
The wave-map formulation and homogeneous geodesics are standard
\cite{ringstrom,lim}. The observation map above singles out a
specific target-coordinate projection.

\section{An exact nonlinear threshold}
Put $s=\log(t/a)$, $\ell=\log(b/a)>0$ and let
$g=P_s(a)$, $q=e^{P(a)}Q_s(a)$ and
$\sigma=(g^2+q^2)^{1/2}$.
Here $Q_s(a)=a Q_t(a)$ is known from the record.
On homogeneous data, \eqref{wave} becomes geodesic motion in $s$.
The conserved speed is $\sigma$, and $e^{2P}Q_s$ is constant.
It follows that $u(s)=e^{P(s)-P(a)}$ obeys
$u_{ss}=\sigma^2u$ with $(u(0),u_s(0))=(1,g)$. Thus the measured
endpoint value is
\begin{equation}
 e^{P(b)-P(a)}
 = F_{q,\ell}(g)
 =\cosh(\sigma\ell)+\frac{g}{\sigma}\sinh(\sigma\ell).
 \label{hom}
\end{equation}
The stationary case is read by continuity.

\begin{theorem}
\label{th:threshold}
Let $h_*>0$ solve $h_*\tanh(h_*/2)=2$.
Consider all homogeneous Gowdy data whose hyperbolic speed is at most $V$.
The observation is injective on this entire class, with $q$ among
the known data, if and only if $V\ell\le h_*$.
If $V\ell>h_*$, some fixed observed $q$ has two distinct compatible
velocities. If $V\ell<h_*$ the homogeneous inverse
is uniformly Lipschitz in $P(b)$ at fixed $q$, with a constant
depending on $V$ and $\ell$.
With no speed bound and $q\ne0$, each observed level strictly above
the minimum has exactly two preimages. The branches meet at a
nondegenerate minimum.
\end{theorem}

\begin{proof}
Nonvertical homogeneous orbits can be parametrized by
$e^{-P}=R\sech r$, $Q=Q_c\pm R\tanh r$ with $r_s=\sigma$.
This is an upper-half-plane semicircle.
Let the initial and terminal parameters be $r_a,r_b$ with
$r_b-r_a=h=\sigma\ell$. The components of $e_P=\partial_P$
in the parallel tangent and normal frame can be chosen to be
$\tanh r$ and $\sech r$.
A velocity variation $\delta g$ at fixed $q$ has zero initial
displacement. Its Jacobi multipliers along and across the target
geodesic are $\ell$ and $\sinh h/\sigma$. Hence
\[
 \frac{\partial P(b)}{\partial g}
 = \frac{\sech r_a\sech r_b}{\sigma}
    \left[h\sinh r_a\sinh r_b+\sinh h\right].
\]
The product of the two hyperbolic sines has minimum
$-\sinh^2(h/2)$, at $r_a=-h/2$. The lower factor in brackets
equals $\sinh h\,[1-(h/2)\tanh(h/2)]$.
For $q=0$ the map is $P(a)+g\ell$, including the stationary case.
Below the threshold the derivative is strictly positive on the
compact speed disc, including the vertical limit by the analytic
expression \eqref{hom}. At equality the derivative is positive
on every nontrivial open admissible interval and may vanish at
its boundary. This proves uniqueness up to and including equality
and the uniform estimate below it.

If $V\ell>h_*$, pick $\sigma<V$ with $\sigma\ell>h_*$ and
choose the semicircle parameter $r_a=-h/2$.
At that velocity the derivative is negative.
Keep $q$ fixed and let $g$ increase to zero, where the derivative
is positive. Both velocities are inside the speed bound, and the
interval between them lies in that bound. The map has a minimum
strictly below both end values on this interval, giving a collision.

For the last assertion, set
$M=\left(\begin{smallmatrix}g&q\\q&-g\end{smallmatrix}\right)$.
Then $F_{q,\ell}(g)=(e^{\ell M})_{11}$.
In the time-ordered expansion about the diagonal matrix
$\operatorname{diag}(g,-g)$, only even numbers of off-diagonal
insertions contribute to this entry. Each term is a positive
integral of $q^{2m}e^{g z}$ with $|z|\le\ell$.
The zero-insertion term is $e^{g\ell}$.
Termwise differentiation is justified by uniform convergence
on compact $g$ intervals, and gives $F''(g)>0$.
For fixed $q\ne0$, \eqref{hom} tends to infinity at both ends
of the real $g$ axis. It has one minimum with positive
second derivative. The logarithm has the same two branches
and a positive second derivative at its minimum.
\end{proof}

For example, take $P(a)=Q(a)=0$, $\ell=3$ and
$\sigma_1=0.95$. Set $q=\sigma_1\sech(3\sigma_1/2)$
and $g_1=-\sigma_1\tanh(3\sigma_1/2)$.
Then $P(b)=0$ by symmetry and the derivative at $g_1$ is negative.
Since $F_{q,3}(0)>1$, there is a second admissible velocity
$g_2\in(g_1,0)$ with $P(b)=0$.
Both speeds are below one. The solution satisfying $g_2$ has
speed about $0.66629$. With window length two the homogeneous
stability bound applies to the entire speed-one class.
The same window already admits a collision when the speed
bound is increased above $h_*/2$, as shown by the example in
Table~\ref{tab:collisions}.

All homogeneous data in this proof satisfy the periodic momentum
constraint. Their metric constants can give identical anchors at $a$.
The two speeds of a collision differ, since
$F_{q,\ell}(g)-F_{q,\ell}(-g)=2g\sinh(\sigma\ell)/\sigma$
for nonzero $g$. In particular their metric derivatives differ.
If one adds a measurement of $\lambda(b)$, then
$\lambda(a)-\lambda(b)=\sigma^2\ell$ gives the speed.
With $\sigma$ known, \eqref{hom} is affine in $g$ and can be inverted.
A second full orbit measurement $Q(b)$ also resolves the ambiguity,
since the initial point and full terminal point determine a unique
target geodesic traversed in time $\ell$.

\begin{table}[ht]
\centering
\begin{tabular}{ccrrrr}
\toprule
$\ell$ & $q$ & $g_1$ & $g_2$ & $\sigma_1$ & $\sigma_2$\\
\midrule
2 & 0.637644 & $-1.357722$ & $-0.693842$ & 1.500000 & 0.942341\\
3 & 0.431979 & $-0.846105$ & $-0.507290$ & 0.950000 & 0.666294\\
\bottomrule
\end{tabular}
\caption{Two collisions with $P(a)=P(b)=Q(a)=0$ in each row.
At each gap both branches use the same initial $q$ and metric anchor.}
\label{tab:collisions}
\end{table}

\input{exact_inhomogeneous.tex}
\input{rank_two.tex}
\subsection{Strictly subcritical collisions in the circular-loop family}
\label{sec:circular-loop}
A distinct exact construction avoids velocity-one sites on
the marked three-torus. It uses the known circular-loop
Gowdy solutions, locally of Bianchi type $\mathrm{VII}_0$
\cite[Section~3]{rendall-twisted} \cite{grubisic}.
The information comparison and the nonlinear shift of the
nodal placement follow directly from their radius equation.

Use hyperbolic geodesic-polar coordinates at a target point
$U_0$. For integer $k\ge1$ write the two maps as
\[
 U_\pm(t,\theta)=\exp_{U_0}\!\left(
 R(t)[\pm\cos(k\theta)e_P+\sin(k\theta)e_Q]\right).
\]
Here the two basis vectors are evaluated at $U_0$,
and a signed radius makes the maps smooth at zero.
The circular-loop reduction gives
\begin{equation}
 R_{tt}+t^{-1}R_t+k^2\sinh R\cosh R=0.
 \label{eq:circular-ode}
\end{equation}
The spatial and temporal target derivatives are orthogonal.
The metric potential is independent of $\theta$ and
the same on both branches with a common anchor.
The target differential has rank two where
$R R_t\ne0$.

\begin{proposition}
\label{prop:circular-loop}
Fix any areal-time ratio $\mathcal R>1$.
For arbitrarily small prescribed nonzero logarithmic
initial radius velocity, there are times $a,b$ with
$b/a=\mathcal R$ and a circular-loop pair with
identical induced spatial metrics and identical $Q_t$
at both times and with identical $\lambda$ throughout.
Their initial $P_t$ differ in the marked data.
The shared asymptotic speed is strictly between zero
and one and is spatially constant.
If $\nu_0=ka$ is a simple linear Dirichlet frequency
for the ratio $\mathcal R$, the small-amplitude
placement has the expansion
\[
 \nu(\epsilon)=\nu_0-
       \frac{\nu_0}{3}
        \frac{\int_1^{\mathcal R} r\,u(r)^4\,dr}
             {\int_1^{\mathcal R} r\,u(r)^2\,dr}
              \epsilon^2+O(\epsilon^4),
\]
where $u$ has initial derivative one at $r=1$ and
vanishes at both ends.
\end{proposition}

\begin{proof}
In the coordinate $r=t/a$ prescribe
$R(1)=0$, $R'(1)=\epsilon$ and set $\nu=ka$.
The equation is
$-(rR')'=\nu^2 r\sinh R\cosh R$.
The quotient $R/\epsilon$ extends analytically and
evenly in $\epsilon$ near zero and follows the
linear radial Dirichlet equation at zero amplitude.
For a linear eigenfrequency $\nu_0$, the derivative
of the terminal value with respect to $\nu$ is
nonzero. Green's identity expresses its product
with $\mathcal R u'(\mathcal R)$ as
$2\nu_0\int_1^{\mathcal R}r u^2\,dr$.
The implicit function theorem gives the nearby
exact placement as an even analytic function of
$\epsilon$. Substitute
$R=\epsilon u+\epsilon^3 R_3+O(\epsilon^5)$,
$\nu=\nu_0+\epsilon^2\nu_2+O(\epsilon^4)$.
Projecting the cubic equation against $u$
and using $\sinh R\cosh R=R+\frac23 R^3+O(R^5)$
gives the stated $\nu_2$.

At the nodal times both maps return to $U_0$.
Their $Q_t$ equals $e^{-P_0}R_t\sin(k\theta)$
on those slices, and their $P_t$ are
$\pm R_t\cos(k\theta)$.
The equal metric potentials follow from equal
wave-map energy sources and the zero momentum source.

For the limit as $t\downarrow0$, set
$x=\log(a/t)\ge0$. The radius solves
$R_{xx}+\nu^2e^{-2x}\sinh R\cosh R=0$ and
\[
 \frac{d}{dx}\left[\frac12 R_x^2+
                \frac12\nu^2e^{-2x}\sinh^2R\right]
 =-\nu^2e^{-2x}\sinh^2R .
\]
Hence $|R_x|\le|\epsilon|$ and
$|R|\le|\epsilon|x$ for all $x\ge0$.
For $|\epsilon|<1$ the restoring force is integrable
there, so $R_x$ has a limit.
The limit divided by $\epsilon$ tends to
$-J_0(\nu_0)$ as $\epsilon\to0$: this is the
normalized linear Bessel limit, and the force
integral converges under the uniform bound just
obtained and dominated convergence.
One can choose a linear Dirichlet frequency with
$J_0(\nu_0)\ne0$ for any fixed ratio.
Indeed $\nu_j=j\pi/(\mathcal R-1)+O(j^{-1})$,
whereas the positive zeros of $J_0$ have expansion
$(m-\tfrac14)\pi+O(m^{-1})$.
If every sufficiently large $\nu_j$ were such a zero,
successive differences would force
$1/(\mathcal R-1)$ to be an integer, contradicting
the quarter-phase shift. Thus the speed is nonzero
for small enough nonzero amplitude and strictly
below one, with the same value at every angle.
\end{proof}

In this circular-loop example a reflection of the
quotient circle exchanges the two histories.
The collision is stated for the marked record.
The separate two-mode Ernst construction supplies
non-equivalence under boundary-preserving
torus-equivariant isometries when quotient
reflections are also allowed.

\subsection{Identical finite metrics with different limiting velocity profiles}
\label{sec:different-velocity-same-metrics}
In the polarized class a common-node pair can also change
the limiting speed map itself while preserving both
induced metrics. In contrast to the rank-two example of
Section~\ref{sec:same-speed-different-modulus}, the
metric ambiguity here includes the leading local speed.

Fix regular times $a,b$ and let
$\mathcal K=\{k\ge1:u_k(b)=0\}$, using
\eqref{eq:invisible}. Consider polarized solutions
with prescribed constant $P(a),P(b)$ and $Q=Q_0$,
with $\lambda(a)$ constant fixed.
Then every compatible $P$ has the form
\begin{equation}
 P(t,\theta)=P(a)+\mu\log(t/a)+
 \sum_{k\in\mathcal K}u_k(t)
       (A_k\cos(k\theta)+B_k\sin(k\theta)),
 \quad \mu=\frac{P(b)-P(a)}{\log(b/a)}.
 \label{eq:polarized-fibre}
\end{equation}
This is the Fourier solution of the scalar equation
with the two prescribed constant displacements.
At both nodes the momentum constraint makes
$\lambda(a),\lambda(b)$ spatially constant.
The mean time quadrature gives
\begin{equation}
 \lambda(a)-\lambda(b)
  =\mu^2\log(b/a)+
     \frac12\sum_{k\in\mathcal K}(A_k^2+B_k^2)
         \int_a^b t\big(u_k'(t)^2+k^2u_k(t)^2\big)\,dt .
 \label{eq:fibre-energy}
\end{equation}
All the weights in the sum are positive.
This describes the full-metric fibre for
these constant observed orbit fields, with no
mode hidden in the zero-frequency part.
In particular, adding the terminal metric
fixes one weighted quadratic radius.

\begin{proposition}
\label{prop:different-velocity-same-metrics}
There exist polarized vacuum three-torus pairs with
identical induced spatial metrics and $Q_t$ at two
regular slices, but different angular profiles of
asymptotic speed. Both profiles can be bounded
strictly between zero and one everywhere.
The profiles can be attained on non-equivalent
histories under torus-equivariant isometries
preserving the observed boundary slices.
The construction is available for a dense set of
areal-time ratios $b/a>1$.
\end{proposition}

\begin{proof}
Choose two distinct modes in $\mathcal K$ with
$J_0(k a)\ne0$ for both. The modes $1,6$
from Proposition~\ref{prop:same-speed-different-modulus}
are one choice.
Use just one mode in each branch of
\eqref{eq:polarized-fibre}, matching the two
positive energies in \eqref{eq:fibre-energy}
by rescaling their nonzero amplitudes.
Here both $Q_t$ traces vanish since $Q$ is
constant. For $0<\mu<1$ and sufficiently
small amplitudes the two limiting profiles are
$\mu+\epsilon_j J_0(k_j a)\cos(k_j\theta)$
with values in $(0,1)$.
The distinct Fourier frequencies make them different.
For a geometric distinction already at the initial
boundary, its quotient metric is constant, so a
torus-equivariant isometry induces a rotation or
reflection of its circle. The mean curvature
depends there on $P_t^2$ with a common nonzero
factor and common additive term, since the
lapse and $\lambda(a)$ agree.
With nonzero $\mu$, the smallest nonzero harmonic
of this quantity is $k_j$ in branch $j$,
which no circle isometry can change.

For density, the ratio $\nu_5(\mathcal R)/
\nu_1(\mathcal R)$ of annular Dirichlet frequencies
is real analytic and nonconstant on $(1,\infty)$.
The limits at its endpoints, $5$ and
$j_{0,5}/j_{0,1}$, differ.
On each open interval of $\mathcal R$ its image
therefore contains an interval.
Choose a rational value $k_2/k_1$ in that image
and put $a=\nu_1(\mathcal R)/k_1$.
The two modes have common nodes.
The values $J_0(\nu_1(\mathcal R))$ and
$J_0(\nu_5(\mathcal R))$ vanish at most on
discrete sets of $\mathcal R$: each eigenvalue
is analytic and nonconstant.
Restrict the interval first to avoid those
discrete sets. This gives the required dense
collection of placements.
\end{proof}

\subsection{Equal finite metrics and asymptotic speed with different limiting moduli}
\label{sec:same-speed-different-modulus}

The exact rank-two pair also separates asymptotic speed
from marked metric scale. Both branches have the same
limiting speed function. Nevertheless a suitable two-mode
phase gives different circle means of the limiting
metric potential even when both induced metrics
were equal at the measured times.

Let $w(\tau,\theta)=A(\theta)\tau+B(\theta)+o(1)$ be
a finite sum of Dirichlet-node seed modes, and let
$\sigma$ be the common speed of the two seed geodesics
with orthonormal initial tangent components $(\pm g,q)$.
Assume $g>0$ and $q\ne0$.
The seed metric potential is identical on both histories.
Write
\[
 D=\log\frac{1+g/\sigma}{1-g/\sigma}.
\]
The explicit hyperbolic geodesics give, wherever
$A(\theta)\ne0$,
\[
 \lim_{\tau\to\infty}
     (\widetilde P_+-\widetilde P_-)
      =-D\operatorname{sgn} A(\theta),\qquad
 \lim_{\tau\to\infty}
     (\widetilde\lambda_+-\widetilde\lambda_-)
      =-2D\operatorname{sgn} A(\theta).
\]
Indeed $e^{P_\pm-P_0}
=\cosh(\sigma w)\pm(g/\sigma)\sinh(\sigma w)$,
and the Ernst transformation changes the metric
potential difference by the same factor two.
The transformed limiting speed of each branch is
$1-\sigma|A(\theta)|$ for sufficiently small amplitude.
At a zero of $A$ the position difference instead tends
to its finite value evaluated at $w=B(\theta)$.
At a simple zero $\theta_0$, the limit with
$\theta=\theta_0+\xi/\tau$ is the corresponding
explicit geodesic ratio evaluated at
$B(\theta_0)+A'(\theta_0)\xi$.
This gives the angular scale of the marked profile
transition. Wherever $A\ne0$, the ratio of the marked
$x$-cycle lengths tends to
$\exp(-D\operatorname{sgn} A/2)$.

\begin{proposition}
\label{prop:same-speed-different-modulus}
There are smooth vacuum three-torus Gowdy pairs with
rank-two target differential on open sets, identical
induced metrics and $Q_t$ at the two measured slices,
and identical asymptotic speed at every angle, whose
limiting metric potentials have different circle means.
\end{proposition}

\begin{proof}
Choose a radial annulus ratio $R$ such that the fifth
Dirichlet frequency is six times the first.
For the rescaled outer radius one, let $z_j=R\nu_j(R)$.
Interval brackets for the first and fifth roots of
$J_0(z/R)Y_0(z)-Y_0(z/R)J_0(z)$ are
\[
\begin{array}{c|cc}
 R & z_1 & z_5\\ \hline
64 & (2.850,2.852) & (15.737,15.739)\\
10^{10} & (2.473,2.475) & (15.007,15.009).
\end{array}
\]
These are verified by the power series for $J_0,Y_0$
with interval arithmetic and geometric bounds on the
series tails. The indices follow from
$j_{0,j}\le z_j\le j\pi/(1-1/R)$:
the lower bound is comparison with the disk by zero
extension in the hole and the upper bound follows
from the transformed radial operator
$-\partial_r^2-1/(4r^2)$.
Continuity gives the required ratio at some $R>64$.
Set $a=z_1/R$ and $b=z_1$, so the angular modes
$1$ and $6$ have common nodes. The upper bound on $z_1$
gives $6a<0.30$, and both relevant $J_0$ values at the
inner radius are positive.

We can then use the two scalar modes with coefficients
chosen to give
\[
 A(\theta)=A_0\big(\cos\theta+\eta\cos(6\theta+\phi)\big)
\]
for small $A_0>0$, small nonzero $\eta$, and
$\phi/\pi$ irrational with $\cos\phi\ne0$.
Indeed the logarithmic slope of each mode is a nonzero
multiple of $J_0(ka)$. The sign-positive set has
length $\pi-2\eta\cos\phi+O(\eta^2)$ by
implicit differentiation at the two zeros
$\pi/2$ and $3\pi/2$. Thus
$\overline{\operatorname{sgn}A}\ne0$.
The initial and terminal metrics and velocities
obey Theorem~\ref{th:rank-two}; the same rank
argument works for the two-mode phase.

To pass to the circle mean, the finite-Bessel seed
expansion is uniform and its metric remainder is bounded.
Also $P_\pm(w)-\sigma|w|$ is uniformly bounded,
so the transformed potentials after subtracting their
common leading $\tau(1-\sigma|A|)^2$ have pointwise
limits and uniform domination. Their circle-mean
difference is therefore
$-2D\,\overline{\operatorname{sgn}A}\ne0$.
\end{proof}

The ratio argument concerns the marked regular-slice
geometry and the displayed asymptotic observables.
It does not infer a change of global spacetime fate
from the velocity-one sites or from the coordinate scale
of the transition.


\section{Linearized inhomogeneous recovery}
For a variation along a homogeneous background, take components
$X,Y$ parallel and perpendicular to the target geodesic.
With Fourier frequency $k\ge1$, covariant linearization of
\eqref{wave} gives
\[
 X''+t^{-1}X'+k^2 X=0,\qquad
 Y''+t^{-1}Y'+(k^2-\sigma^2/t^2)Y=0.
\]
The transverse sign also follows by representing the geodesic
as a vertical target line with $P=\sigma\log t+$constant.
The linearized $Q$ equation there has damping $(1+2\sigma)Q_t/t$.
Its target-frame coordinate $Y=e^P\delta Q$ has the displayed
negative potential. With the rescaling by $\sqrt t$ the two
potentials become $\beta_+/t^2$ and $\beta_-/t^2$, where
$\beta_+=1/4$, $\beta_-=1/4-\sigma^2$.

Let $S_\beta(b;k)$ be the value at $b$ of the solution of
$z''+(k^2+\beta/t^2)z=0$ with $z(a)=0$, $z'(a)=k$.
Hold $P(a),Q(a),Q_t(a)$ fixed and vary $P_t(a)$.
The single terminal observation $P(b)$ then has multiplier
$\sqrt{a/b}\,T_k$ when the unknown is $\delta P_t(a)/k$, where
\begin{equation}
 T_k=A S_{\beta_+}(b;k)+B S_{\beta_-}(b;k),\qquad
 A=\tanh r_a\tanh r_b,\quad B=\sech r_a\sech r_b .
 \label{mix}
\end{equation}
Both the starting direction and the terminal projection are
fixed by the observation protocol. Two Jacobi solutions are involved
even though the measured terminal quantity is a scalar.

\begin{theorem}
\label{th:fourier}
Take $b-a=\pi p/q_*$ with $p,q_*$ positive coprime integers.
There is a subcritical background of speed $1/\sqrt2$
for which $|T_k|\asymp k^{-3}$ on the large multiples of $q_*$,
whereas it is bounded away from zero on all large nonmultiples.
Outside a locally finite exceptional set of reference times
the full linearized inverse exists. Its loss is exactly
three derivatives for $\delta P_t(a)/k$.
Adding $Q(b)$ or $\lambda(b)$ to the observation restores the
one-derivative loss for this background, with the same exclusion
of exceptional reference times.
\end{theorem}

\begin{proof}
Choose the semicircle with $\sinh r_a\sinh r_b=1$.
It exists because this product, as the center in log time varies,
ranges from $-\sinh^2((r_b-r_a)/2)$ to infinity.
Then $A=B>0$ and the two potentials have coefficients $\pm1/4$.
In the standard Hankel expansion, after extracting
$\sqrt{2/(\pi z)}e^{i(z-\nu\pi/2-\pi/4)}$ with
$\nu^2=1/4-\beta$, the prefactor is
\[
 H_\beta(z)
 =1-\frac{i\beta}{2z}-\frac{\beta(\beta+2)}{8z^2}
       +\frac{i\beta(\beta+2)(\beta+6)}{48z^3}+O(z^{-4})
\]
\cite[\S 10.17]{dlmf}. The Bessel Wronskian gives
$S_\beta(b;k)=\operatorname{Im}
[e^{ik(b-a)}\overline{H_\beta(ka)}H_\beta(kb)]$,
where the exact prefactor is used in this identity.
At resonance, the order-$k^{-1}$ coefficient is
$\beta(u-v)/2$ up to parity, with $u=1/a,v=1/b$.
It cancels in \eqref{mix}. The even-in-$\beta$ coefficient
of order $k^{-3}$ is
\begin{equation}
 -\beta^2(u-v)\left(\frac{u^2+v^2}{6}+\frac{uv}{24}\right),
 \label{coefficient}
\end{equation}
also up to parity. It is nonzero. The order-$k^{-2}$ term is real
before multiplication by $e^{ik(b-a)}$ and vanishes at resonance.
For nonmultiples, the leading term is $2B\sin(k(b-a))$.
These estimates are uniform for $a$ on compact positive subintervals.

For completeness, each fixed nonzero mode has only isolated zeros
as a function of $a$. The transfer is analytic in $a>0$ at fixed
gap $d=b-a$. On a resonant mode, averaging the two Volterra
series cancels all odd orders in the potential. As $a\to\infty$,
the second-order term has leading value equal to that of a constant
potential $\beta/a^2$, with an error $O(a^{-5})$, since
$(a+x)^{-2}=a^{-2}+O(a^{-3})$ uniformly for $0\le x\le d$.
Expanding $k\sin(d\sqrt{k^2+W})/\sqrt{k^2+W}$ for constant $W$
gives second-order coefficient $(-1)^{kp/q_*}(-3d/(8k^3))$.
Thus the averaged term is nonzero for sufficiently large $a$.
Nonresonant modes approach the nonzero free multiplier as
$a\to\infty$. Uniform high-frequency estimates then imply local
finiteness of the exceptional set. The zero spatial mode has
positive observation derivative on these backgrounds by the
same-sign Jacobi weights in the preceding proof.

Adding $Q(b)$ reads both target-frame components separately.
Adding $\lambda(b)$ reads the tangent component instead:
linearizing the metric quadrature along the geodesic gives
$\delta\lambda(b)=-2\sigma X(b)$ for a fixed initial anchor.
The tangential input component is nonzero here.
At resonance the individual nonzero channels are of order
$k^{-1}$ with opposite signs before projection, so either added
measurement avoids the cancellation.
Away from resonance, the original scalar observation has a
uniform high-frequency lower bound. Finite modes use the
nonvanishing hypothesis already established.
\end{proof}

For $a=\pi/4$ and $b=7\pi/16$, the predicted limit of
$k^3|T_k|$ on the resonant subsequence is $0.01302349475$.
High-precision Bessel evaluations at $k=16384$ give $0.01302349439$.
The paired orbit measurement has $k$ times its normalized
multiplier norm approaching $0.06820926$.
Independent integrations of the coordinate linearization checked
the trace formulas on modes $1,15,16,32,49$.
The derivative loss above refers to the normalized velocity.
Recovering the unscaled $P_t$ costs one additional derivative.

\input{detuning.tex}

\section{An exact linearized kernel and a cubic nonlinear separation}
For $\sigma=1/2$ place $r_a=0$. The initial variation in $P_t$ lies
entirely in the normal Jacobi channel and its effective potential
after rescaling vanishes. Thus
$T_k=\sech(r_b)\sin(k(b-a))$.
At $k(b-a)\in\pi\mathbb{Z}$, both terminal orbit variations vanish.
The first-order metric variation vanishes as well because
$X=0$ throughout the interval. This example lies within the
speed-one class and can satisfy the homogeneous positive
threshold. The kernel involves inhomogeneous velocity
perturbations and is absent from the purely homogeneous inverse.

\begin{proposition}
For any such resonant mode, the full nonlinear map that reads
the two induced spatial metrics, along with the prescribed
$Q_t(a)$, cannot have a pairwise local inverse with H\"older
exponent exceeding $1/3$ for the missing $P_t(a)$.
The same obstruction holds for the original smaller record.
On the chosen signed Fourier perturbations the terminal $P$
discrepancy has nonzero cubic order for sufficiently large
reference $a$ at fixed resonant gap.
\end{proposition}

\begin{proof}
Choose smooth initial data with identical $P(a),Q(a),Q_t(a)$
and with $P_t(a)$ perturbed by $\pm\epsilon\cos(k\theta)$
about the apex background.
The initial position is homogeneous in $\theta$, so the periodic
momentum constraint is satisfied and both initial metric
anchors can be fixed identically.
Regular-time smooth dependence on this scalar parameter follows from
the semilinear wave system on the compact time interval,
by differentiating its evolution equations and applying local
energy estimates about the smooth background.
The difference between the two observations has zero first-order
coefficient by the kernel calculation. Its even terms cancel.
It is $O(\epsilon^3)$ in any fixed Sobolev norm of the smooth
output family. The missing initial velocities differ by order
$\epsilon$. An inverse bound with exponent greater than $1/3$
is therefore impossible for arbitrary pairs in such a neighborhood.

To see the cubic mechanism, use Fermi coordinates $N,L$
with target metric $dN^2+\cosh^2N\,dL^2$. On the background,
$N_0=0$ and $L_0'=\sigma/t$.
Reflection $N\mapsto-N$ makes $L$ even and $N$ odd in $\epsilon$.
Write $N=\epsilon N_1+\epsilon^3N_3+O(\epsilon^5)$ and
$L=L_0+\epsilon^2 L_2+O(\epsilon^4)$.
The equation for $L_2$ is driven by $-2 L_0' N_1 N_{1,t}$.
The cubic forcing in the normal equation is
$2L_0' L_{2,t}N_1+(2/3)(L_0')^2N_1^3$.
There is no initial cubic term because the fixed initial target
point has an exactly linear velocity perturbation in the
Fermi normal direction.

Let $x=t-a$ and $d=b-a=m\pi/k$.
For $a\to\infty$, the leading normal coefficient is
$u(x)\cos(k\theta)$ with $u(x)=\sin(kx)/k$.
On $[0,d]$, the mode-zero and mode-$2k$ time derivatives
of $a L_2$ tend to
$-\sigma u^2/2$ and $-\sigma x\sin(2kx)/(4k)$, respectively,
by their forced constant-coefficient equations with zero initial data.
The mode-$k$ forcing for $a^2N_3$ therefore tends to
\[
 -\frac{\sigma^2}{2}u^3
 -\frac{\sigma^2}{4k}xu\sin(2kx).
\]
The Green function at $d$ is $(-1)^{m+1}u(x)$.
Using $u^2\sin(2kx)=(k/2)(u^4)'$ and $u(0)=u(d)=0$
gives the nonzero terminal limit
\[
 \lim_{a\to\infty} a^2(N_3)_k(b)
 =(-1)^m\frac{9\sigma^2 d}{64 k^4}.
\]
The limit follows by continuous dependence for the finitely
many forced harmonic equations for these coefficients.
At the terminal slice $N_1=0$, so the odd cubic term of
$P$ is the projection of $N_3$ onto $dP$, with factor
$\sech(r_b)>0$. Thus this family has a genuine cubic separation
for sufficiently large $a$.
\end{proof}

An independent numerical test used $a=1$, $b=1+\pi/2$ and $k=2$.
A Fourier evolution of the full coordinate equations was compared
with an ODE for the expansion coefficients in Fermi coordinates.
At amplitude $0.025$, the $\cos(2\theta)$ coefficient of the terminal
$P$ difference divided by $\epsilon^3$ was $-0.000736959$.
The Fermi expansion gave $-0.000736964$. The combined terminal
$P,Q,\lambda$ difference decreased by factors near eight when
the signed perturbation amplitude was halved.
This illustrates the cubic mechanism on regular finite slices.

\section{Stable recovery from three fixed slices}
\label{sec:rank-two-coercivity}
Frequency-assigned sampling gives a general nonlinear recovery
result. Parallel transport from the two characteristic directions
can separate their contributions in a full orbit snapshot.
This yields local nonlinear stability with a fixed number of
slices near a broad class of inhomogeneous backgrounds,
including the explicit Gowdy-to-Ernst backgrounds of
Section~\ref{sec:rank-two}.

Let $\bar U=(\bar P,\bar Q)$ be a smooth background, use
$t=a+d$, and let $V$ be a linearized wave map with
$V(a)=0$, $D_tV(a)=f e_P$. The smooth orthonormal frame
$e_P=\partial_P$, $e_Q=e^{-\bar P}\partial_Q$ has connection
one-form $\omega=e^{\bar P}d\bar Q$, with
$D_\mu e_P=\omega_\mu e_Q$.
Let $A_\pm(t,\theta)$ be the frame component vector obtained
by transporting $e_P$ along the $\pm$ null rays from
$(a,\theta\mp d)$ to $(t,\theta)$.
For $h=\partial_\theta^{-1}(f-\langle f\rangle)$ of mean zero,
the leading displacement in this frame is
\begin{equation}
 \sqrt{t/a}\,V(t,\theta)
  =\frac12\left[A_-(t,\theta)h(\theta+d)
                    -A_+(t,\theta)h(\theta-d)\right]+R_t f .
 \label{eq:transport-leading}
\end{equation}
For integer $s\ge1$, $R_t$ maps $H^{s-1}$ boundedly to
$H^{s+1}$ on any smooth regular-time interval in question.

To verify the remainder estimate, multiply the covariant
linearization by $\sqrt{t/a}$ to remove $t^{-1}D_t$.
Its principal part is $D_t^2-D_\theta^2$ and its other
terms are a smooth zeroth-order potential given by curvature
and $1/(4t^2)$. Apply this operator to the leading expression
in \eqref{eq:transport-leading}. The terms with $h''$
cancel between equal-speed time and space derivatives.
The terms with $h'$ vanish by the two transport equations
$(D_t\pm D_\theta)A_\pm=0$.
Only smooth coefficients times $h$ remain.
At $a$ the leading expression is zero and its initial
derivative equals $f-\langle f\rangle$ in the $e_P$ direction
up to a smooth coefficient times $h$.
The error therefore has zero initial displacement, initial
velocity bounded in $H^s$ by $\|f\|_{H^{s-1}}$, and source
with that same regularity. The linear energy estimate proves
the claim. The formula is insensitive to the finite-rank
choice of treatment of the mean.


\subsection{A criterion from the relative null transport}
At a chosen terminal time $c=a+d$, let $\alpha_\pm(\theta)$
be the angles of $A_\pm$, defined by the signed integrals
$-\int_{\gamma_\pm}\omega$ on the two null rays.
Write $\rho=(\alpha_- -\alpha_+)/2$, $\beta=\sin\rho$,
and rotate the output frame by the mean angle
$(\alpha_-+\alpha_+)/2$.
The two leading components in \eqref{eq:transport-leading},
up to the common factor $1/2$, then read
\[
 \cos\rho\,\big(H-\mathcal T H\big),\qquad
 \beta\,\big(H+\mathcal T H\big),
 \quad H(\theta)=h(\theta+d),\quad
 \mathcal T H(\theta)=H(\theta-2d).
\]
We use the conditions
\begin{equation}
 \inf_\theta|\cos\rho(\theta)|>0,\qquad
 \inf_{\theta\in\mathbb T}\sum_{j=0}^{J}
           |\beta(\theta-2jd)|^2>0
 \quad\text{for some finite }J .
 \label{eq:rotation-cover}
\end{equation}
For an irrational $d/\pi$ and $\beta\not\equiv0$,
the second condition follows from density of
circle-rotation orbits and a finite subcover of the
translates of an open set where $|\beta|$ is bounded below.

In particular the conditions hold at all sufficiently
short delays after a slice where $Q_\theta(a)$ is
nonzero somewhere. Indeed, expansion of the two ray
integrals about $(a,\theta)$ gives
\begin{equation}
 2\rho=2d\,\omega_\theta(a,\theta)
        +d^2\big(\partial_t\omega_\theta-
                  \partial_\theta\omega_t\big)(a,\theta)
        +O(d^3).
 \label{eq:relative-expansion}
\end{equation}
If $Q_\theta(a)\equiv0$ but $P_\theta(a)Q_t(a)$
is nonzero somewhere, the first term vanishes and the
coefficient of $d^2$ is nonzero on an interval since
$d\omega=e^{P}dP\wedge dQ$. These two cases include any
background of rank two at some point of the initial slice.
In either case a fixed open arc has $|\beta|$ bounded
below by a positive multiple of $d$ or $d^2$ for small
$d$. If the arc is wider than $2d$, the first
$\lceil\pi/d\rceil+1$ translates cover the circle.
The short-delay result therefore holds for rational
as well as irrational placements.

For the explicit Gowdy-to-Ernst construction,
$\omega_t=C e^{-P(w)}w_\theta$ and
$\omega_\theta=C e^{-P(w)}w_t$.
Set $B'(z)=C e^{-P(z)}$ with $B(0)=0$.
The two rotations are $-B(w(c,\theta))$ and
$+B(w(c,\theta))$, since the seed phase vanishes at $a$.
For a small-amplitude seed of a single mode $k$,
a time $c$ with nonzero mode amplitude and
$2kd\notin\pi\mathbb Z$ satisfies both conditions,
with $J=1$.

\begin{theorem}
\label{th:fixed-rank-two}
Consider a smooth Gowdy background and a terminal time
$c$ satisfying \eqref{eq:rotation-cover}. In addition to the initial
$P(a),Q(a),Q_t(a)$ and the metric anchor, read $P,Q$ at $c$
and finitely many low Fourier coefficients of $P$ at a
suitably chosen $a+d_1\in(a,c)$.
There is an $H^s$ neighborhood of the background initial
$P_t$ for each integer $s\ge1$ in which the nonlinear
observation is locally bi-Lipschitz from the missing
velocity in $H^s$ to these observations, using $H^{s+1}$
for the full orbit fields and a fixed finite-dimensional
norm for the early coefficients. The same estimate applies on the
periodic momentum-constraint subspace.
\end{theorem}

\begin{proof}
The displayed principal components control
$H-\mathcal T H$ and $\beta H$ in $L^2$.
If $b_j(\theta)=\beta(\theta-2jd)$, then
\[
 b_j H = \mathcal T^j(\beta H)+b_j(H-\mathcal T^j H).
\]
Telescoping bounds the second term by
$j\|\beta\|_\infty\|H-\mathcal T H\|_{L^2}$.
Summing these estimates and using \eqref{eq:rotation-cover}
controls $\|H\|_{L^2}$.
Apply the same estimate to $\partial_\theta^{s+1}H$.
Commutation through the smooth multiplication coefficients
adds only terms bounded by $\|H\|_{H^s}$.
Add the compact remainder from \eqref{eq:transport-leading}
and the mean mode. For the exact derivative
$\mathcal A_c:f\mapsto (\delta P(c),\delta Q(c))$,
this gives
\begin{equation}
 \|f\|_{H^s}\le C
    \big(\|\mathcal A_c f\|_{H^{s+1}}+\|f\|_{H^{s-1}}\big).
 \label{eq:rank-two-compact}
\end{equation}
The frame factors and the scalar time rescaling
are bounded invertible multipliers on the output space.


The low-frequency reading can be chosen constructively.
Let $C$ be the constant in \eqref{eq:rank-two-compact},
and choose $K$ so that $C/\langle K\rangle<1/4$.
The $H^{s-1}$ norm splits into the low part and a tail
bounded by $\langle K\rangle^{-1}\|f\|_{H^s}$.
For $d_1<\min(d,1/K)$, apply the free-wave Duhamel formula
to the coordinate linearization from $a$ to $a+d_1$.
The map $f\mapsto\Pi_K\delta P(a+d_1)/d_1$
differs from the free diagonal map
$\operatorname{sinc}(k d_1)\Pi_Kf$
by a map bounded from $H^{s-1}$ to $H^{s-1}$ with norm
$O(d_1)$. This follows from the smooth coefficient
energy estimate and $|S_k(r)|\le r$.
The free gains are at least $\sin(1)$.
Choose $d_1$ smaller if needed to absorb
the $O(d_1)\|f\|_{H^{s-1}}$ contribution into the
left-hand side of \eqref{eq:rank-two-compact}.
The joint derivative is bounded below.
Equivalently, compactness in \eqref{eq:rank-two-compact}
shows that the full-terminal kernel is finite dimensional,
and the early reading removes it.

For the nonlinear step, the first-order energy formulation
with coordinate source $\mathcal F$ gives a $C^1$ solution
map on a neighborhood of the smooth background, on the fixed
time interval. One can obtain the derivative by subtracting the
linearized equation from the difference equation and applying
the same energy estimate to the quadratic Taylor remainder.
The derivative is continuous in operator norm by a linear
energy estimate on the difference of two coefficient fields.
Thus on a sufficiently small convex ball the observation
derivative differs from its value at the background by
less than half the lower bound of the joint derivative.
Integrating this operator difference along the segment between
two velocities proves the pairwise lower Lipschitz estimate.
The upper estimate follows from the local forward energy bound.
The affine periodic momentum constraint can be imposed after
the estimate.
\end{proof}

This result applies to full orbit observations at a time
where the two null transports differ.
At a common node of the seed phase, $\rho=0$ everywhere
and the argument supplies only the translation difference.
It does not give a stability estimate at the resonant
two-metric pair. The construction uses two later time
slots, and the earlier slot needs only finitely many
low Fourier coefficients. It leaves open
whether that extra reading can be omitted or replaced
by a particular previously imposed terminal time.

\subsection{When the full terminal orbit snapshot suffices}

\begin{proposition}
\label{prop:two-slice-open}
There is an open class of smooth Cauchy backgrounds,
including solutions with rank-two target differential on
open sets, near which the nonlinear missing-velocity
inverse is locally bi-Lipschitz from one full terminal
orbit snapshot and the initial $P(a),Q(a),Q_t(a)$.
The norms are $H^s$ for velocity and $H^{s+1}$ for the
snapshot, for integer $s\ge1$.
One explicit choice starts at a geodesic-image background
$U=\gamma(w)$ with $\gamma$ a unit-speed nonvertical
geodesic. Suppose $w(a,\theta)$ is nonconstant and its
range lies in a compact segment on which
$|\langle e_P,\dot\gamma\rangle|>0$.
A sufficiently short delay $d=\pi/L$ with $L$ an integer
gives the conclusion.
\end{proposition}

\begin{proof}
On a geodesic-image background the tangential and normal
linearized equations decouple in the frame parallel along
$\gamma$. The curvature terms take values in the normal
direction. Thus the tangent component $X$ obeys
$X_{tt}+t^{-1}X_t-X_{\theta\theta}=0$ with initial
position zero and velocity
$X_t(a)=\langle e_P,\dot\gamma(w(a))\rangle f$.
The full terminal orbit snapshot contains $X(c)$.
The scalar rational-gap bound of
Proposition~\ref{prop:safe-gap} gives, for sufficiently
large $L$,
\[
 \|X_t(a)\|_{H^{s-1}}\le C\|X(c)\|_{H^{s+1}} .
\]
The multiplier $\langle e_P,\dot\gamma(w(a))\rangle$
is smoothly invertible by hypothesis. In addition,
$Q_\theta(a)$ is nonzero somewhere by nonverticality
and nonconstant $w(a)$. The same short delay satisfies
\eqref{eq:rotation-cover}. We may therefore use the
tangent estimate to control the $H^{s-1}$ remainder
in \eqref{eq:rank-two-compact} and obtain a bounded-below
derivative for the single terminal full orbit snapshot.
The $C^1$ argument of Theorem~\ref{th:fixed-rank-two}
turns this into a local nonlinear inverse. Continuity
in the background Cauchy data preserves the strictly
positive lower bound at the same time. A small initial
velocity perturbation perpendicular to the geodesic
preserves the momentum constraint and creates rank two
where $w_\theta(a)\ne0$. This gives genuinely rank-two
members of the open class.
\end{proof}

\paragraph{Necessity of the cover at principal order.}
When $\cos\rho$ is bounded away from zero,
\eqref{eq:rotation-cover} is also necessary for the
one-terminal linearized orbit observation to have an
energy lower bound, before any lower-frequency
obstruction is considered. Indeed, if the cover fails,
compactness gives a point whose full forward rotation
orbit lies in the zero set of $\beta$.
For irrational rotation this forces $\beta\equiv0$,
and the translation difference has arbitrarily small
Fourier gains by rational approximation.
For rational rotation the zero orbit is finite.
Choose smooth, rotation-invariant envelopes supported in
shrinking neighborhoods of that orbit and multiply them
by arbitrarily high Fourier carriers invariant under the
same rotation. The principal translation difference
vanishes and the multiplication by $\beta$ becomes
arbitrarily small. Taking the carrier large for each
fixed envelope makes its contribution from the
lower-order remainder tend to zero in the normalized
energy scale. These sequences contradict the required
lower bound. This argument concerns the full orbit
snapshot without an additional terminal $\lambda$ reading.

\subsection{Generic two-slice stability near inhomogeneous Cauchy data}
\label{sec:generic-two}

\begin{theorem}
\label{th:generic-two}
Fix any smooth periodic vacuum Gowdy Cauchy background
at a regular time $a>0$ with $Q(a,\theta)$ nonconstant.
There is a common short rational delay $d=\pi/L$ and
a smooth analytic path of admissible Cauchy data starting
at this background, keeping $Q(a)$ fixed, such that
the full orbit snapshot at $a+d$ has a locally
bi-Lipschitz missing-$P_t$ inverse at every point of
the path outside a finite exceptional set.
The velocity and snapshot norms are respectively
$H^s$ and $H^{s+1}$ for integer $s\ge1$.
In particular the given background can be approximated
by admissible backgrounds with stable two-slice
recovery, using the same delay. Stability persists
in an open neighborhood of each nonexceptional point.
\end{theorem}

\begin{proof}
Choose a nonvertical target geodesic whose $Q$ range
contains the range of the given $Q(a)$ on a compact
subarc away from its turning point. This is possible by
taking a sufficiently large upper-half-plane semicircle
with its center to one side of the prescribed range.
Set the comparison initial $P$ to the value on that
arc and take both comparison velocities zero.
These data have a geodesic-image evolution.
They meet the hypotheses of Proposition~\ref{prop:two-slice-open}.

Interpolate the initial $P,P_t,Q_t$ linearly in a
parameter $\zeta\in[0,1]$, from the given data to this
comparison datum, keeping $Q(a)$ fixed.
To impose the vacuum periodic momentum constraint,
replace the trial $Q_t$ by
\[
 Q_t-\frac{\int_{\mathbb T}
        (P_\theta P_t+e^{2P}Q_\theta Q_t)\,d\theta}
       {\int_{\mathbb T}e^{2P}Q_\theta^2\,d\theta}\,Q_\theta .
\]
The denominator is strictly positive on the whole
path. The correction vanishes at the endpoints.
All Cauchy data vary real analytically in $\zeta$ in
all fixed smooth Sobolev norms.

There is a common short time interval of forward
existence, including a small real extension of the
parameter path beyond both endpoints. On a fixed
arc $|\omega_\theta(a)|=e^{P(a)}|Q_\theta(a)|$
has a uniform positive lower bound, since $Q(a)$
is fixed. The expansion \eqref{eq:relative-expansion}
therefore gives \eqref{eq:rotation-cover} uniformly
for a sufficiently short $d>0$ on the path.
Choose $d=\pi/L$ small enough for the same conclusion
and the scalar tangent estimate at the comparison datum.

Let $\mathcal A_\zeta:H^s\to (H^{s+1})^2$
be the derivative of the full terminal orbit snapshot
with respect to the missing initial $P_t$.
The estimate \eqref{eq:rank-two-compact} holds for
all real $\zeta$ in a neighborhood of the path.
Thus $\mathcal A_\zeta$ has closed range and finite
kernel, and its Hilbert normal operator
$B_\zeta=\mathcal A_\zeta^\ast\mathcal A_\zeta$
is Fredholm of index zero.
The first-order energy formulation is real analytic
in Cauchy data on a common small time interval.
For instance, one may complexify the energy spaces and
use the analytic implicit function theorem on the
Duhamel equation on short subintervals, since the
coordinate nonlinearity is analytic on the Sobolev
algebra. This also gives analyticity of the observation
derivative in bounded-operator norm.
The normal operators consequently extend to a holomorphic
Fredholm family on an open connected complex
neighborhood of the real parameter path, using
$B_z=\mathcal A_{\bar z}^\ast\mathcal A_z$.
Fredholmness on a possibly thinner neighborhood follows
from its openness in operator norm.

At the comparison endpoint, the tangent argument in
Proposition~\ref{prop:two-slice-open} gives a
bounded-below $\mathcal A_1$; hence $B_1$ is invertible.
The analytic Fredholm theorem makes the real
noninvertibility set discrete in a neighborhood of
$[0,1]$, and therefore finite on that compact interval.
At any remaining point the derivative is bounded below.
Continuity of the derivative and the same pairwise
integration argument as in
Theorem~\ref{th:fixed-rank-two} give the nonlinear
inverse on its Cauchy neighborhood.
\end{proof}

There is no contradiction between density of stable
nearby backgrounds at a short chosen delay and the exact
different histories at the specially resonant delays.
The theorem varies a Cauchy background at a fixed short
delay; it makes no claim about a specified background
and a specified arbitrary observation gap.

\subsection{Replacing the blind terminal slice in the exact rank-two pair}
\label{sec:exact-rank-recovery}

\begin{theorem}
\label{th:exact-rank-recovery}
Take the pair of Theorem~\ref{th:rank-two}
with $gq\ne0$ and a single-mode seed at nodes $a<b$.
Write its phase as $w_\epsilon=\epsilon u_{k_0}(t)\cos(k_0\theta)$.
There is an explicit rational delay $d=\pi/L$ with
$0<d<b-a$ such that, for all sufficiently small
nonzero $|\epsilon|$, the full orbit fields at
$c=a+d$ give a locally bi-Lipschitz missing-$P_t(a)$
inverse around either rank-two history.
For each integer $s\ge1$, the snapshot is measured
in $H^{s+1}$ and the velocity in $H^s$.
The lower bound for the derivative can be taken
to be a positive constant times $|\epsilon|$.
\end{theorem}

\begin{proof}
Choose an integer $L>2k_0$ satisfying the short-gap
conditions of Proposition~\ref{prop:safe-gap}, and
large enough that $u_{k_0}(c)\ne0$.
Let $\mathcal P_\epsilon$ and $\mathcal Q_\epsilon$
denote the two linearized terminal observation
operators in the scaled orthonormal frame of
\eqref{eq:transport-leading}. Their dependence on
$\epsilon$ in the operator norm $H^s\to H^{s+1}$
is analytic on a small interval.
At $\epsilon=0$ the transformed background is a
vertical homogeneous solution with $Q$ constant.
The $P$-variation then solves the scalar
Euler--Poisson--Darboux equation and the $Q$-variation
vanishes for the given missing-velocity direction.
Moreover, the first derivative of the $P$-observation
in $\epsilon$ vanishes. In the coordinate linearized
$P$ equation the terms beyond the free damped wave
are proportional to $Q_t^2-Q_\theta^2$ times
$\delta P$, or to a first derivative of $Q$ times
a first derivative of $\delta Q$.
Both are quadratic in the amplitude at this
reference. The time rescaling of $P$ is independent
of $\epsilon$. Hence
\[
 \mathcal P_\epsilon=\mathcal P_0+O(\epsilon^2),
 \qquad
 \mathcal Q_\epsilon=\epsilon\mathcal Q_1+O(\epsilon^2)
\]
in the stated operator norms.

On the multiples of $L$, the two shifts of the
zero-mean primitive $h=\partial_\theta^{-1}f$ in
\eqref{eq:transport-leading} agree, with sign
$(-1)^{k/L}$ on mode $k$.
The rotation angle has first derivative
$q\,u_{k_0}(c)\cos(k_0\theta)$ at $\epsilon=0$,
because $B'(0)=q$ for this normalization.
Thus on the nonzero multiples of $L$ the leading
part of $\mathcal Q_1$ is
\[
 f\longmapsto q\,u_{k_0}(c)\cos(k_0\theta)\,
                 \mathcal J\,\partial_\theta^{-1}f ,
\]
where $\mathcal J$ is the sign operator
$(-1)^{k/L}$.
The differentiated remainder is bounded
from $H^{s-1}$ to $H^{s+1}$, as follows by
differentiating the smooth transport coefficients
and the forced remainder equation.
Multiplication by $\cos(k_0\theta)$ sends each
multiple of $L$ to two disjoint sidebands
$k\pm k_0$. They do not overlap with sidebands
of any other multiple of $L$ because $L>2k_0$.
The two shifted Sobolev weights are comparable to
the unshifted ones. The displayed leading operator
therefore has a uniform positive lower bound.
After discarding finitely many resonant modes,
the compact remainder has arbitrarily small norm
on the remaining high resonant subspace.
Consequently $\mathcal Q_1$ is bounded below there.

Split $f=f_G+f_B$ into the complementary modes
(the nonmultiples and finitely many small
multiples of $L$) and the remaining large
multiples. The scalar short-gap transfer
$\mathcal P_0$ is bounded below on $G$ by the
nonresonant free phases and the nonzero finite
transfers of Proposition~\ref{prop:safe-gap}.
Projecting the $P$ record onto $G$ thus gives
\[
 \|f_G\|_{H^s}\le C\big(
    \|\mathcal P_\epsilon f\|_{H^{s+1}}+
      \epsilon^2\|f\|_{H^s}\big).
\]
The first-order estimate on $B$, together with
the triangle inequality for $\mathcal Q_1 f_G$,
gives
\[
 \|f_B\|_{H^s}\le C'\left(
   \frac{\|\mathcal Q_\epsilon f\|_{H^{s+1}}}{|\epsilon|}
       +\|f_G\|_{H^s}+|\epsilon|\|f\|_{H^s}\right).
\]
Combining and absorbing the small-amplitude terms
bounds the inverse by $O(|\epsilon|^{-1})$.
The constants may depend on the selected times
and the seed parameters.
The two sign branches use the same argument.
The local nonlinear inverse follows by continuity
of the Cauchy derivative as in
Theorem~\ref{th:fixed-rank-two}.
\end{proof}

The order of the lower gain is sharp as
$\epsilon\to0$ for the rational placement used here.
One can take arbitrarily high multiples of $L$.
The scalar transfer of $\mathcal P_0$ then tends
to zero in energy normalization, while the
uniform remainder estimates give contributions of
order $\epsilon^2$ and $|\epsilon|$ respectively
in the $P$ and $Q$ channels.

\subsection{Which extra terminal field supplies stable information}
\begin{proposition}
\label{prop:rank-two-fields}
For the exact small-amplitude rank-two pair in
Theorem~\ref{th:exact-rank-recovery}, the off-node
observation $P(c),\lambda(c)$ gives the same local
energy-norm stability even without a terminal $Q(c)$.
It is enough to read $\lambda_\theta(c)$ in $H^s$.
In contrast, no fixed finite collection of later $P$
snapshots alone gives an energy-normalized local
Lipschitz inverse at these backgrounds.
\end{proposition}

\begin{proof}
Write the two transport vectors in the frame as
$A_\pm=(\cos\rho,\mp\sin\rho)$.
The momentum constraint gives the leading
mean-zero part of the rescaled metric variation as
\[
 -\big(C_-H-C_+\mathcal T H\big),\qquad
 C_-=\langle A_-,U_t+U_\theta\rangle,\quad
 C_+=\langle A_+,U_t-U_\theta\rangle .
\]
Here the rescaling of $\delta\lambda$ is $\sqrt{c/a}/c$.
This follows by substituting the two leading orbit
transport terms in the momentum constraint and taking
a spatial primitive, as in Section~\ref{sec:metric-positive}.
All omitted terms map $H^{s-1}$ to $H^{s+1}$.
In the Ernst family the half difference of the
coefficients is
\[
 \frac{C_--C_+}{2}
 =\cos\rho\,\widetilde P_\theta
    +\sin\rho\,e^{\widetilde P}\widetilde Q_t
 =w_\theta[-\cos\rho\,P_w+\sin\rho\,C e^{-P(w)}]
 =-g w_\theta .
\]
The last equality holds because the seed geodesic
tangent with initial components $(g,q)$ is parallel:
its components at $w$ are obtained by a rotation
through $-B(w)=-\rho$.
At zero amplitude the linearized
$\lambda_\theta(c)$ equals $2\delta P_\theta(c)$.
Subtract this known channel and take the rescaled
mean-zero primitive of the difference.
Its linearization in amplitude, on modes that
are multiples of $L$, has the principal term
\[
 -2g k_0u_{k_0}(c)\sin(k_0\theta)\,
               \mathcal J\,\partial_\theta^{-1}f .
\]
This is bounded below on the high resonant
subspace by the same disjoint-sideband argument,
since $g\ne0$ and $L>2k_0$.
The remainder and the good-mode $P$ bounds are
the same as in Theorem~\ref{th:exact-rank-recovery}.
Thus $P$ and the metric channel give a bounded-below
derivative, and then a local nonlinear inverse.

For the negative statement, at each later time $t_i$
the leading $P$ component is a smooth multiplier
$\cos\rho_i$ times the difference of the two translated
primitives. The shifts are opposites and the $P$
components of their transported vectors are equal.
By simultaneous rational approximation choose a
sequence $k_m\to\infty$ with
$\max_i|\sin(k_m(t_i-a))|\to0$.
The corresponding unit-velocity Fourier modes in $H^s$
give leading $P$ norms tending to zero at every
chosen time. The compact remainders tend to zero
as well by \eqref{eq:transport-leading}.
A pairwise local Lipschitz inverse for these $P$
observations would imply a lower bound for the
derivative at the reference background, a
contradiction.
\end{proof}

\subsection{Observing the terminal metric potential}
\label{sec:metric-positive}
The spatial derivative of the remaining metric potential
measures an additional component of the orbit velocity.
This produces a positive result even at polarized
inhomogeneous backgrounds, where the two orbit transports
coincide.

For the linearized displacement of
\eqref{eq:transport-leading} put $X=\sqrt{c/a}\,V(c)$.
Linearizing the momentum constraint at $c$ gives
\[
 -\frac{\sqrt{c/a}}{2c}\,\delta\lambda_\theta(c)
  =\langle D_t X-X/(2c),U_\theta\rangle
   +\langle D_\theta X,U_t\rangle ,
\]
where $X$ in its time derivative denotes the full
time-rescaled displacement. The measured terminal orbit
fields determine the terms containing only $X$ and
its spatial derivative. Thus an observation of
$\lambda_\theta(c)$ gives $\langle D_tX,U_\theta\rangle$.
Its mean-zero spatial primitive has principal term
\[
 \tfrac12\left[\langle U_\theta,A_-\rangle H+
                    \langle U_\theta,A_+\rangle\mathcal T H\right].
\]
Indeed, the derivative of the two shifted primitives
in \eqref{eq:transport-leading} has the sum sign,
and commuting the spatial primitive through a smooth
coefficient gains an additional derivative.
The mean terms are finite rank.
The same forced-wave energy estimate used in
\eqref{eq:transport-leading} gives this principal
formula modulo a map from $H^{s-1}$ to $H^{s+1}$.

With $e_b$ the unit bisector from the mean rotation,
define $b_2=\cos\rho\,\langle U_\theta,e_b\rangle$.
Subtract the already observed translation difference
from the last principal expression. This controls
$b_2H$ in addition to $\beta H$.
Consequently the proof of
\eqref{eq:rank-two-compact} and
Theorem~\ref{th:fixed-rank-two} applies with
\[
 \inf_\theta\sum_{j=0}^{J}
  \big(|\beta(\theta-2jd)|^2+|b_2(\theta-2jd)|^2\big)>0
\]
and with $\cos\rho$ bounded away from zero.
For sufficiently short delays this holds if either
$P_\theta(a)$ or $Q_\theta(a)$ is nonzero somewhere.
For the first condition, $b_2$ converges uniformly to
$P_\theta(a)$ as $c\downarrow a$.
For the second, use \eqref{eq:relative-expansion}.
In either case the nonzero interval is covered by
finitely many translations. The resulting local estimate
uses the full terminal induced metric, plus a finite
low-mode $P$ reading if a lower-frequency kernel remains.

\begin{proposition}
\label{prop:two-metric-polarized}
Near a polarized background with nonconstant $P(a)$,
one sufficiently short rational terminal placement gives
a locally bi-Lipschitz missing-$P_t$ inverse from the
full terminal induced metric and the original initial
data. The same conclusion holds on an open class of
nearby Cauchy backgrounds. The observation norm includes
$\lambda_\theta(c)$ in $H^s$ in addition to
$P(c),Q(c)$ in $H^{s+1}$.
\end{proposition}

\begin{proof}
At the polarized reference, the $P$-variation solves
the scalar Euler--Poisson--Darboux equation and the $Q$
variation vanishes for the specified perturbation.
Apply the one-derivative scalar estimate at a short
rational gap as in Proposition~\ref{prop:safe-gap}
to bound the $H^{s-1}$ norm of the missing velocity by
the terminal $H^{s+1}$ norm of $P$.
Use it to remove the lower-order term of the
full-metric estimate above. Continuity of the
derivative and integration along a small neighborhood
give the local nonlinear statement.
For continuity on the unconstrained velocity space
one may use the periodic constraint source
$-2t\langle U_t,U_\theta\rangle$ as the output for all
reduced solutions. On periodic vacuum solutions it is
exactly $\lambda_\theta$. Restriction to initial
velocities compatible with the momentum constraint
therefore gives the claimed geometric observation.
\end{proof}

The analytic-path argument of Theorem~\ref{th:generic-two}
also applies to this full-metric output when $P(a)$ is
nonconstant. Keep $P(a)$ fixed, interpolate to a
polarized comparison state, and enforce the constraint
by subtracting a multiple of $P_\theta(a)$ from $P_t(a)$.
The geometric cover remains uniform along the path
because the same nonzero arc of $P_\theta(a)$ persists.
This gives nearby admissible backgrounds with stable
two-slice full-metric recovery at a common short delay
when either initial orbit field is spatially nonconstant.

\section{Recovering both initial velocities from orbit snapshots}
\label{sec:both-velocities}
The original record supplies one component of the
initial velocity. When both are withheld, the two
null-ray transports act on a vector rather than on
one initial direction. A closed sequence of ray shifts
then gives a matrix condition for stable recovery.
This section concerns two full orbit snapshots and one
initial metric anchor. It does not give a negative
result when a full metric potential is supplied as an
additional spatial field.

Let the missing velocity variation be
$f\in H^s(\mathbb T;\mathbb R^2)$ in the chosen
orthonormal frame. The initial orbit position is fixed.
Replace the columns $A_\pm$ of
\eqref{eq:transport-leading} by the full parallel
transport matrices on the two null rays, and let $h$
be the componentwise mean-zero primitive of $f$.
The same leading formula and compact remainder hold.
After rotating the terminal output by $A_-^{-1}$,
the principal operator on $H(\theta)=h(\theta+d)$ is
\[
 I-\mathcal U_d,\qquad
 (\mathcal U_d H)(\theta)
  =W_d(\theta)H(\theta-2d),\qquad
 W_d=A_-^{-1}A_+\in \mathrm{SO}(2).
\]
This is a unitary weighted rotation in $L^2$.
An equivalent $H^{s+1}$ conjugation changes the
operator by a compact term because the coefficient
matrix is smooth.

\begin{proposition}[A rational ray-holonomy test]
\label{prop:both-ray-holonomy}
For irrational $d/\pi$ the linearized two-orbit-snapshot
map cannot have an energy-norm lower bound for both
missing velocity components, even after excluding
finitely many directions.
For rational $d/\pi=p/q$ in lowest terms define
\[
 \mathcal H_d(\theta)=
     \prod_{j=0}^{q-1} W_d(\theta-2jd) .
\]
The principal operator is bounded below precisely when
$I-\mathcal H_d$ is uniformly invertible.
In that case the exact observation derivative has a
lower bound up to an $H^{s-1}$ velocity remainder.
A sufficiently early finite set of extra low-mode
orbit observations gives a bounded-below derivative
and a local nonlinear Lipschitz inverse.
\end{proposition}

\begin{proof}
For irrational rotation, conjugating $\mathcal U_d$
by multiplication by $e^{im\theta}$ on the
complexified space multiplies it by $e^{-2imd}$.
Its nonempty essential spectrum is therefore the
whole unit circle. Since the operator is unitary,
the point $1$ has a weakly null approximate sequence.
Conjugation to the Sobolev norm and the smoothing
observation remainder do not remove this sequence.
The mean and the linearized integral momentum
constraint amount to finitely many restrictions,
also harmless for the weakly null sequence.

For rational rotation $\mathcal U_d^q$ is
multiplication by $\mathcal H_d$.
If $I-\mathcal H_d$ is uniformly invertible,
factor $I-\mathcal U_d^q$ to obtain the claimed
principal lower bound:
\[
 I-\mathcal U_d^q
    =(I+\mathcal U_d+\cdots+\mathcal U_d^{q-1})(I-\mathcal U_d).
\]
If $\mathcal H_d(\theta_0)$ has eigenvalue one, use
an envelope supported near $\theta_0$ and a vector
fixed by that holonomy. Sum its first $q$ weighted
translates; they have disjoint supports for small
enough envelopes, and the residual telescopes to
$(I-\mathcal U_d^q)$ applied to the original envelope.
It becomes arbitrarily small as the support shrinks.
Attach a high Fourier carrier with frequency a
multiple of $q$ to commute with each shift.
Taking account of the Sobolev normalization and
the compact remainders by choosing arbitrarily
large carriers shows failure of an energy lower
bound even after compact corrections.
The vector energy estimate and the earlier low-mode
argument of Theorem~\ref{th:fixed-rank-two} prove
the final claim. The local nonlinear argument
is unchanged when both initial velocities vary
subject to the affine momentum constraint.
\end{proof}

\begin{theorem}
\label{th:both-circular}
For a fixed sufficiently short rational delay,
there are circular-loop Gowdy backgrounds, and open
neighborhoods containing general nearby Cauchy data,
on which two full orbit snapshots are locally
bi-Lipschitz for both missing initial velocities.
The construction works for all sufficiently small
nonzero initial radii outside a countable set.
\end{theorem}

\begin{proof}
Take a circular-loop solution with angular winding
$k_0\ge1$, initial radius $\rho>0$, and zero initial
radial velocity. The chosen interval is short
enough that the radius stays positive for small
$\rho$.
In the radial and unit angular frame,
the spatial connection is $k_0\cosh R(t)$ times
the generator of rotations, and the temporal
connection vanishes. Hence the angle of the
holonomy in Proposition~\ref{prop:both-ray-holonomy}
is
\[
 -2k_0 q\int_a^c\cosh R(t)\,dt
   \equiv -2k_0q\int_a^c(\cosh R(t)-1)\,dt
                   \pmod{2\pi}.
\]
For any sufficiently small $\rho>0$ the
last integral is positive and the angle
is strictly between $-2\pi$ and zero.
The principal lower bound therefore holds.

In this frame all coefficients of the exact
linearized equation are independent of $\theta$,
so there is a $2\times2$ terminal transfer matrix
on each complex Fourier mode.
These matrices depend analytically on $\rho$,
including $\rho=0$ by using the rotating
orthonormal frame at the target center.
At $\rho=0$ the equation is the constant-background
scalar Gowdy equation on each physical component.
A Fourier mode $k$ written in the rotating frame
has the two shifts $k\pm k_0$ in the fixed frame.
Choose the rational delay as in
Proposition~\ref{prop:safe-gap}, so every one of
those scalar transfers, including zero frequency,
is nonzero. For each fixed $k$ the transfer
determinant is analytic and nonzero at
$\rho=0$. The union of its zero sets over all
integer $k$ is countable.
Away from this set the derivative has no kernel.
The lower bound up to compact remainder then
gives a full energy lower bound, by the same
compactness argument as before.
Continuity of the derivative in the Cauchy
data proves stability in nearby fully nonlinear
solutions.
\end{proof}


\paragraph{The holonomy of the Ernst pair.}
At a short rational delay $d=\pi/L$ with $L>2k_0$,
the rank-two Ernst pair also satisfies the principal
test for both velocities when its amplitude is
small and nonzero.
Its relative ray rotation is $-2B(w(c,\theta))$.
Use $B(w)=q w-(qg/2)w^2+O(w^3)$.
Summing the first harmonic over the $L$ ray shifts
gives zero, while the sum of its square is $L/2$.
The resulting holonomy angle is therefore
\[
 \frac{qgL}{2}\,\epsilon^2u_{k_0}(c)^2+O(\epsilon^3),
\]
uniformly in $\theta$ and with the sign of the
chosen geodesic branch. It is nonzero for small
nonzero amplitude and $\;gq\ne0$.
One full terminal orbit snapshot consequently
gives the high-frequency estimate in this family.
A finite set of early low-mode coefficients
removes any finite-dimensional kernel.
The next theorem gives explicit two-terminal
placements that do not require locating this kernel.


\paragraph{The local two-time compatibility equation.}
At the circular-loop backgrounds of
Theorem~\ref{th:both-circular}, the bounded-below
two-component derivative is also onto the two
terminal orbit fields. To see this, extend the
mean-zero primitive to an isomorphism from $H^s$
to $H^{s+1}$ by adding the mean mode.
Its difference from the leading expression is
finite rank. The rational holonomy condition
makes its composition with $I-\mathcal U_d$
invertible. The exact observation derivative is
a compact perturbation of this invertible
operator, hence has index zero.
Its trivial kernel then gives surjectivity.

It follows from the inverse function theorem that
the reduced wave-map problem with prescribed orbit
positions on the two slices has a locally smooth
solution for initial positions and terminal values
in their natural Sobolev neighborhoods.
The periodic vacuum metric exists with a specified
anchor exactly on the image of the momentum condition
\[
 \int_{\mathbb T}\langle U_t(a),U_\theta(a)\rangle\,d\theta=0 .
\]
When $U_\theta(a)$ is nonzero somewhere, this is
locally one compatibility equation on the measured
orbit positions. Indeed its derivative in the recovered
velocity is a nonzero linear functional, and the
two-snapshot reconstruction is a local diffeomorphism
for the reduced equations. At the circular-loop
reference with the initial orbit position held fixed,
the linearized compatibility condition is that the
mean angular component of the terminal displacement
in the rotating frame vanish. The initial angular
velocity enters the constraint only through its
zero Fourier mode, and that mode is decoupled
with nonzero scalar transfer in this frame.

\subsection{Three orbit snapshots in the exact rank-two family}

\begin{theorem}
\label{th:both-ernst}
At each sufficiently small nonzero amplitude of the
single-mode rank-two Gowdy-to-Ernst pair, two suitably
placed later full orbit snapshots together with the
initial full orbit snapshot recover \emph{both} missing
initial velocities locally with a Lipschitz estimate
in the wave-energy Sobolev norms. The placements
can be chosen as different short rational gaps,
independently of the amplitude.
\end{theorem}

\begin{proof}
Let the two delays be $d_i=\pi/L_i$, $i=1,2$,
with both integers $L_i$ sufficiently large,
distinct, and greater than $2k_0$.
At the zero-amplitude vertical background, the two
time-scaled linearized channels have inverse-square
potentials $1/(4t^2)$ and $-3/(4t^2)$.
At these short rational placements neither scalar
channel has a finite-frequency kernel.
For the negative potential the Pr\"ufer correction
to $kd_i$ has the opposite sign from the
positive-potential case and absolute magnitude
less than $\pi/L_i$ for large $k$.
For remaining modes choose $L_i$ large so that
$\varphi(c_i)<\pi$; the phase remains strictly
positive since its derivative is $k$ at each
zero of $\varphi$. The zero spatial mode has
nonzero Jacobi transfers.
On high frequencies both channels have leading
normalized gain $\sin(kd_i)$.

Write $\mathcal A_\epsilon$ for the joint derivative
with respect to the unknown two-component velocity
and $\epsilon$ for the seed amplitude.
It has the norm expansion
$\mathcal A_\epsilon=\mathcal A_0+
\epsilon \mathcal A_1+O(\epsilon^2)$.
Let $L$ be the least common multiple of $L_1,L_2$.
Split the velocity into $f_B$, its high modes
divisible by $L$, and $f_G$, the remaining modes.
The free joint map $\mathcal A_0$ is bounded below on
$G$, since all finite transfers are nonzero and
each large noncommon multiple has a nonzero free
gain at one of the times.

At each output Fourier mode, project the four
observations orthogonally away from the two
columns of $\mathcal A_0$ there, in the
energy-normalized coordinates. Call this bounded
projection $\Pi^\perp$.
For an input high common multiple $k$, the first
variation of the transport angles maps to sidebands
$k\pm k_0$. Up to the common rotation of the
velocity vector and Fourier normalization, its
two time coefficients at each sideband are
\[
 \big((-1)^{k/L_1}u_{k_0}(a+d_1),
      (-1)^{k/L_2}u_{k_0}(a+d_2)\big).
\]
Indeed the two shifted primitives coincide at each
common resonance and the rotations change with
opposite signs on the two rays.
The leading free column at that sideband has time
coefficients proportional to
\[
 \big((-1)^{k/L_1}\sin(k_0d_1),
      (-1)^{k/L_2}\sin(k_0d_2)\big).
\]
The two vectors are not collinear: near zero,
\[
 \frac{u_{k_0}(a+d)}{\sin(k_0d)}
    =\frac1{a k_0}-\frac{d}{2a^2 k_0}+O(d^2)
\]
is strictly decreasing. The relative signs at
the two sidebands do not affect the projected gain.
Distinct high common multiples give disjoint
sidebands because $L>2k_0$.
The exact free projections at those sidebands
approach the displayed limiting projection with
error $O(1/|k|)$, uniformly along the resonant
subsequence. The differentiated transport remainder
is compact as before. Consequently
$\Pi^\perp\mathcal A_1$ is bounded below on $B$
after increasing the finite cutoff if needed.

The good-mode estimate and the projected estimate
therefore give
\begin{align*}
 \|f_G\|_{H^s}
   &\le C\big(\|\mathcal A_\epsilon f\|+
                 |\epsilon|\|f\|_{H^s}\big),\\
 \|f_B\|_{H^s}
   &\le C'\big(\|\mathcal A_\epsilon f\|/|\epsilon|+
                  \|f_G\|_{H^s}+|\epsilon|\|f\|_{H^s}\big).
\end{align*}
For small nonzero amplitude the error terms
are absorbed. The exact derivative is then bounded
below, and continuity of the nonlinear derivative
gives the local result.
\end{proof}

\section{The metric can remove irrational-time weak modes}
\label{sec:full-metric-holonomy}

The irrational-time obstruction in
Proposition~\ref{prop:both-ray-holonomy}
concerns observations of the orbit metric alone.
The remaining spatial metric coefficient gives
a scalar reading of the otherwise hidden
transported velocity. This reading has a
different principal symbol from the two orbit
coordinates.

Use the notation of Section~\ref{sec:both-velocities}.
In addition to both orbit fields at $c$,
measure the metric gradient $\lambda_\theta(c)$
in $H^s$, equivalently the mean-zero metric
potential in $H^{s+1}$. Use the initial orbit
fields and a metric anchor at $a$.
After the full orbit observation is used,
the principal metric reading is proportional to
\[
 B(\theta)\cdot H(\theta),\qquad
 B(\theta)=A_-^{-1}(c,\theta)U_\theta(c,\theta).
\]
In fact the derivative of the momentum constraint,
after taking a spatial primitive, gives to principal order
\[
 -\Big(\langle A_-H,U_t+U_\theta\rangle
       -\langle A_+\mathcal T H,U_t-U_\theta\rangle\Big).
\]
The contribution involving $U_t-U_\theta$ paired
with $A_-H-A_+\mathcal T H$ is already observed.
The remaining contribution is $-2 B\cdot H$.
The common scale of this expression is $\sqrt{ac}$.
Differentiating the transported primitive gives
this expression by the momentum equation, whose
linearization is
$-2t(\langle D_t V,U_\theta\rangle+
\langle U_t,D_\theta V\rangle)$.
The differentiated transport coefficients and
the forced error in \eqref{eq:transport-leading}
give terms bounded from $H^{s-1}$ to $H^{s+1}$
after the primitive. The mean is a finite-rank term.

\begin{theorem}[The metric projection along successive rays]
\label{th:full-metric-gramian}
Let $d/\pi$ be irrational and let
$B_j(\theta)=(\mathcal U_d^j B)(\theta)$.
The linearized observation of the terminal
orbit fields and metric gradient is bounded below
in the $H^s$ velocity norm up to an
$H^{s-1}$ remainder if and only if
for some finite $J$ and $b>0$,
\begin{equation}
 \sum_{j=0}^J B_j(\theta)B_j(\theta)^\mathsf{T}\ \ge bI
          \quad\text{for every }\theta .
 \label{eq:metric-gramian}
\end{equation}
The sufficient direction holds at rational delays
as well. Under \eqref{eq:metric-gramian},
finitely many early orbit coefficients remove
any finite-dimensional kernel, giving a local
nonlinear Lipschitz inverse for both velocities.
For sufficiently short irrational delays, the
criterion holds if $U_\theta(a)$ and
$D_\theta U_\theta(a)$ are independent at
some angle of the initial orbit curve. In particular
this applies at the common initial node of the
Ernst pair, where the orbit curve is a nonconstant
horocycle arc in the upper half-plane.
\end{theorem}

\begin{proof}
Write $\mathcal U=\mathcal U_d$.
Orthogonality of the transport gives the identity
\[
 B_j\cdot H
    =\mathcal T^j(B\cdot H)
         +B_j\cdot(H-\mathcal U^j H).
\]
The second term is controlled by telescoping
$H-\mathcal U H$ and boundedness of $B_j$.
Condition \eqref{eq:metric-gramian} then controls
$\|H\|_{L^2}$ from the two readings.
Applying the same argument to the top spatial
derivative gives the Sobolev estimate up to
one lower order. The smoothing terms in the
linearized wave and metric maps give the
claimed $H^{s-1}$ remainder.
Reading the early low orbit modes and taking
the nonlinear neighborhood
uses the same proof as
Theorem~\ref{th:fixed-rank-two}, with two
components of initial velocity in the energy.

Conversely, if \eqref{eq:metric-gramian}
fails for every $J$, for arbitrarily long finite
orbit segments there is a unit vector at one end
orthogonal to all transported $B$ readings on
that segment. Transport this vector back to the
first angle of the segment. Place an $L^2$-unit
smooth envelope $h_0$ around that angle, pointing
in the transported vector direction. The supports
of $\mathcal U^j h_0$ for $0\le j\le J$ can be
made disjoint by irrationality.
Their $B$ readings are uniformly as small as
desired by continuity and the orthogonality at
the centers, after choosing the envelope narrow
enough for this finite segment.
The normalized sum
$(J+1)^{-1/2}\sum_{j=0}^J\mathcal U^j h_0$
has unit norm and its difference under $I-\mathcal U$
has norm at most $2/\sqrt{J+1}$.
Multiplication by a high Fourier carrier with
$e^{-2imd}$ arbitrarily near $1$ preserves these
bounds up to an arbitrarily small error.
Taking the carriers sufficiently high yields a
weakly null sequence in $L^2$.
Conjugating to $H^{s+1}$ and imposing the
finite mean and momentum restrictions introduce
only compact or finite-rank corrections, which
vanish in this sequence. The claimed lower
estimate must therefore fail.

To see the short-delay criterion, compare
$B_0$ and $B_1$ in the terminal frame at
$(c,\theta)$ using $A_-(c,\theta)$.
They become $U_\theta(c,\theta)$ and
$U_\theta(c,\theta-2d)$ transported to $(c,\theta)$
across the two rays with common initial point
$(a,\theta-d)$. The first-order difference between
these two vectors is $-2dD_\theta U_\theta(a,\theta)$.
Indeed the spatial displacement is $2d$ and the
first-order temporal contributions of the reversed
and forward transports cancel. Thus the two
vectors are independent on an open arc for
sufficiently small $d$. By irrationality every
rotation orbit enters that arc and finitely many
translates cover the circle. Orthogonal transport
preserves independence, giving
\eqref{eq:metric-gramian} with finitely many
readings in the principal estimate.
\end{proof}

\paragraph{The complete high-frequency test for two full metrics.}
If the initial metric profile is supplied, its gradient
also reads the initial velocity through
$\delta\lambda_\theta(a)=-2a\langle f,U_\theta(a)\rangle$.
After taking the spatial primitive and shifting,
this adds the principal row
$B^{(a)}(\theta)=U_\theta(a,\theta+d)$ in the initial
orthonormal frame of $H$.
The extra differentiated coefficient term is
one order smoother. Denote the available rows by
$B^{(\ell)}$, with $\ell$ running over the terminal
metric row and the initial row if supplied.
For irrational $d/\pi$ the same proof gives the
necessary and sufficient test up to compact
remainders with the sum in
\eqref{eq:metric-gramian} taken over these rows.
Their joint mean readings are finite rank.

For rational $d/\pi=p/q$ the corresponding necessary
and sufficient condition up to compact remainders is
uniform positive definiteness of
\begin{equation}
 (I-\mathcal H_d)^\mathsf{T}(I-\mathcal H_d)
 +\sum_{\ell}\sum_{j=0}^{q-1}
   (\mathcal U_d^jB^{(\ell)})
   (\mathcal U_d^jB^{(\ell)})^\mathsf{T}.
 \label{eq:rational-full-metric}
\end{equation}
Indeed $(I-\mathcal H_d)H=H-\mathcal U_d^qH$
is bounded by $q\|(I-\mathcal U_d)H\|$.
Each projected row is controlled by the telescoping
identity in the proof above. Conversely, if the
displayed matrix is singular at an angle, choose
a unit vector fixed by its closed transport and
orthogonal to every projected row.
Use normalized narrow envelopes about its finite
rotation orbit with the transported vector
directions. Choose their relative signs by the
transport and sum them. The residual in
$I-\mathcal U_d$ is small because the last
transport returns the initial direction at the
center of the envelope. All metric readings are
small by orthogonality at the corresponding centers.
Multiplying by arbitrarily high carriers at
multiples of $q$ commutes with the shift and
gives a weakly null sequence. The compact
lower-order corrections then vanish, proving
necessity. This test includes pure closed-ray
holonomy when no metric row is supplied, and
can still be positive at identity holonomy when
the projected metric directions span the plane.

\begin{proposition}[Two slices at irrational separation]
\label{prop:irrational-metric}
For every positive initial radius in the circular-loop
family with zero initial radial velocity, all sufficiently
short positive delays give a local energy-norm
Lipschitz inverse for both velocities from the two
induced metrics. The delays may be rational or irrational.
One may replace the initial metric by the orbit
part and use a single metric anchor. This stability
persists in a neighborhood of the Cauchy data
and on an open interval of nearby known observation
delays.
\end{proposition}

\begin{proof}
In the rotating orthonormal frame of the circular
loop, $B=A_-^{-1}U_\theta$ is constant in $\theta$
and is nonzero when the final radius is positive.
The matrix $W_d$ is also constant in this frame.
Choose the delay short enough that its rotation
angle $-2k_0\int_a^c\cosh R(t)\,dt$ is strictly
between $-\pi$ and zero.
Then $B$ and $W_d B$ are independent and
\eqref{eq:metric-gramian} holds with $J=1$.
It remains to rule out an exact kernel. The linearized
equations decompose in Fourier modes of the
rotating frame. If the terminal orbit and
metric readings vanish, the terminal orbit variation
vanishes and the gradient constraint gives
$\langle D_t V(c),U_\theta(c)\rangle=0$.
Thus for each Fourier mode the terminal derivative
is radial. Evolve the mode backward from zero
terminal position and radial derivative $g e_R$.
In the reversed time variable, after multiplying by
$\sqrt{t/c}$ the exact linearized operator differs
from the frozen-connection operator by a potential
with norm at most $D_0+D_1|k|$, where
\[
 D_0=\sup_{[a,c]}\left(
 |\Omega^2-\Omega_0^2|+\|K\|+\frac{1}{4t^2}\right),
 \qquad
 D_1=2\sup_{[a,c]}|\Omega-\Omega_0|.
\]
Here $\Omega(t)=k_0\cosh R(t)$, $\Omega_0=\Omega(c)$,
and the curvature potential $K$ in the polar
frame has diagonal entries $(k_0\sinh R)^2,-R_t^2$.
Indeed the principal spatial operator on mode $k$
is the square of $ikI+\Omega J$, with $J$ the
rotation generator. The time rescaling removes
$t^{-1}\partial_t$ and adds $1/(4t^2)$.

For the frozen operator the two circular
polarizations have frequencies $k\pm\Omega_0$.
The radial vector has equal projection magnitude
onto them. Put $q_k=\max(d^{-1},|k|)$
and assume $\Omega_0d<\pi/4$ and $\Omega_0d<1/2$.
The free reverse transfer from zero position
is diagonal with entries $\sin((k\pm\Omega_0)r)/
(k\pm\Omega_0)$, using the continuous value at
zero. Its norm times $q_k$ is at most $2$ for
$0\le r\le d$. At $r=d$ the radial input has
output norm times $q_k$ at least
\[
 \frac{\sin(\Omega_0d)}{1+\Omega_0 d}|g|.
\]
This follows by bounding
$|k\pm\Omega_0|\le(1+\Omega_0d)q_k$ and using the
sum of the two squared sines. The weighted
Duhamel formula bounds the correction by
$2\big(\exp(2d(dD_0+D_1))-1\big)|g|$.
The correction is $O(d^2)$ since $D_0$ is bounded
and $D_1=O(d)$, while the free lower bound is of
order $d$. Thus all sufficiently short delays
force nonzero displacement at $a$ for any nonzero
radial terminal derivative, uniformly over $k$.
This rules out the kernel and also gives a direct
lower estimate. For nonzero terminal readings,
the momentum constraint controls the tangential
terminal derivative by the gradient reading and
$q_k$ times the terminal orbit displacement.
Separating that tangential derivative and the measured
displacement in the reverse problem, the same
Duhamel estimate bounds their contribution at $a$
in the $q_k$-weighted position norm. Since the
position at $a$ is prescribed, the radial lower
bound then controls the remaining terminal
derivative by these readings. The regular-time
energy estimate transfers this bound to the initial
velocity. For example $a=1.6$,
$k_0=4$, $R(a)=0.32$, $R_t(a)=0$ allow every
$0<d\le0.02$. The decreasing radial energy
gives $|R|\le0.32$, $|R_t|<1.31$ and $R>0.29$ on
the interval. In particular $4\le\Omega_0<4.21$,
$D_0<2.11$ and $D_1<3.44d$. These bounds
make the correction less than $0.45d$
and the displayed free gain greater than $3.68d$
over that interval.

The compact-remainder bound for the enlarged
metric observation and injectivity now give the
energy-norm lower bound. The nonlinear estimate
again follows from continuity of the Cauchy derivative.

The positive lower bound in the finite sum
\eqref{eq:metric-gramian} persists continuously
in the delay, along with the bounded lower-order
remainder. Translation is strongly continuous on
Sobolev spaces even though it need not be continuous
in operator norm. If the full lower bound failed
arbitrarily close to an injective delay $d_0$, take
unit $H^s$ velocities $f_n$ at $d_n\to d_0$
with terminal observations tending to zero.
After a subsequence $f_n$ converges in $H^{s-1}$
and weakly in $H^s$ to $f$.
The regular-time energy bound and strong
continuity in that weaker norm imply that $f$
is invisible at $d_0$. Injectivity gives $f=0$,
and the uniform lower bound up to an $H^{s-1}$
remainder contradicts the unit $H^s$ norms.
The Cauchy derivative varies with initial data
in operator norm uniformly on the fixed compact
time interval, so the nonlinear neighborhood
can also be chosen uniformly on a smaller
interval of delays.
\end{proof}

The orbit-only two-slice map still has the irrational
weak modes of Proposition~\ref{prop:both-ray-holonomy}
even on the circular backgrounds in this proposition.
For instance, the two projected directions
for a short circular loop are just the tangent
of the orbit curve and its rotated copy.
This makes the effect of the metric gradient visible
in the same evolution and at the same observation times.

\paragraph{Longer placements can make both full metrics weak.}
The shortness hypothesis has geometric content even
in the circular family. Use the smooth rotating
frame defined by parallel transport along the signed
radial geodesics, including at zero radius.
For any winding, choose the later time so that
\[
  \int_a^c k_0\cosh R(t)\,dt=\pi .
\]
For every small positive initial radius this defines
one delay near $\pi/k_0$, by strict monotonicity
in $c$. It is strictly shorter than $\pi/k_0$ and
depends analytically on the initial radius, so the
delays fill a nontrivial interval and are irrational
relative to $\pi$ outside a countable set of small
radii. To see that the endpoint loop can stay
nontrivial, take $a$ large with $k_0$ fixed.
The normalized linear radial solution with zero
initial derivative then approaches
$\cos(k_0(t-a))$ uniformly on this interval
as $a\to\infty$, by the scalar Duhamel formula
with damping $1/t$. Its terminal value is negative
for sufficiently large $a$ and remains so after
dividing the nearby nonlinear radius by its
initial amplitude.

At the chosen delay $W_d$ is the identity in
the rotating frame and the metric directions
at the two times are both angular, since the
single-ray rotation is a half turn. All their
translates are collinear. For the irrational
choices the two-full-metric criterion therefore
fails. This gives arbitrarily weak energy-normalized
velocity directions at longer exact placements,
even when both full induced metrics are supplied.

\section{From finite observations to a low-velocity asymptotic state}
\label{sec:low-velocity-tail}

A stable missing-velocity inverse can be composed with
a stable evolution toward the singular end in a
uniformly subcritical, small-data region. The following
estimate gives one explicit overlap between those
two restrictions. The asymptotic expansion itself
belongs to the established low-velocity Gowdy
theory \cite{ringstrom,ringstromavtd,libounces}.
Here the estimate controls differences between
limiting states by differences between finite
measurements.

Let $r=\tau-\tau_0\ge0$, fix $0<v_*<1$, and set
\[
 V=P_\tau,\quad W=e^{-\tau}P_\theta,\quad
 X=e^P Q_\tau,\quad Y=e^{P-\tau}Q_\theta,\qquad
 Z=(V-v_*,W,X,Y).
\]
Use the Sobolev norms with weights
$\langle k\rangle^s=(1+k^2)^{s/2}$ and normalized
Fourier coefficients on the $2\pi$ circle.
Let $A_s=2^s(1+2(1+1/(2s-3)))^{1/2}$.
This bounds the algebra and embedding constants at
orders $s$ and $s-1$, by Fourier convolution and
$\sum_{k\ge1}k^{-2(s-1)}\le1+1/(2s-3)$.
Define the momentum $\Pi=e^{2P}Q_\tau$.
The normalization of the limits uses the same $\tau_0$
for all solutions, so the constants in the expansions
are the intercepts relative to $r$.

\begin{proposition}
\label{prop:low-velocity-difference}
Fix $s\ge3$, $\tau_0$ and $v_*$. Set
$\delta=\min(v_*,1-v_*)/(16A_s)$.
Every smooth periodic vacuum datum with
$\|Z(\tau_0)\|_{H^s}<\delta$ and the Gowdy momentum
constraint evolves to $\tau=\infty$. In a bounded
set of the initial $H^s$ orbit fields, the map from
initial orbit data, velocities and one metric anchor
to
\begin{align*}
 (v,p,q)&=\lim_{r\to\infty}(P_\tau,P-rv,Q),\\
 (\lambda_\infty,\Pi_\infty)&=
                 \lim_{r\to\infty}(\lambda-rv^2,\Pi)
\end{align*}
is locally Lipschitz into
$(H^{s-1})^2\times H^s\times H^{s-1}\times H^{s-2}$.
Take the norm of the input difference to be
\[
 \mathcal I=\|Z(\tau_0)-\widehat Z(\tau_0)\|_{H^s}
      +\|(P,Q,\lambda)(\tau_0)-
                   (\widehat P,\widehat Q,\widehat\lambda)(\tau_0)\|_{H^s}.
\]
One may use the initial $H^{s+1}$ orbit fields, $H^s$
velocities and the scalar metric anchor to bound
$\mathcal I$. Here $v$ is the limiting hyperbolic
speed, with $0<v<1$, and
\begin{equation}
 \partial_\theta\lambda_\infty
 =
 2v\,\partial_\theta p+2\Pi_\infty\,\partial_\theta q .
 \label{eq:tail-momentum}
\end{equation}
\end{proposition}

\begin{proof}
The first-order equations are
\begin{align*}
 V_\tau &= e^{-\tau}W_\theta+X^2-Y^2,\\
 W_\tau &= e^{-\tau}V_\theta-W,\\
 X_\tau &= e^{-\tau}Y_\theta-V X+W Y,\\
 Y_\tau &= e^{-\tau}X_\theta+(V-1)Y-W X .
\end{align*}
Use the squared sum of the component $H^s$ norms
for the energy. Each paired principal spatial
term cancels by periodic integration by parts.
For $D=(W,X,Y)$ and $c=\min(v_*,1-v_*)$,
the diagonal terms give damping at least
$c\|D\|_{H^s}^2$ in the half-energy derivative.
Every nonlinear source contribution to it is
bounded in sum by
$3A_s\|Z\|_{H^s}\|D\|_{H^s}^2$, using the Sobolev
algebra inequality. Smallness therefore gives
\begin{equation}
 \|Z(\tau)\|_{H^s}^2+
             c\int_{\tau_0}^{\tau}\|D(\sigma)\|_{H^s}^2\,d\sigma
     \le \|Z(\tau_0)\|_{H^s}^2.
 \label{eq:low-velocity-energy}
\end{equation}
Sobolev embedding also gives
$0<v_*/2<V< (1+v_*)/2<1$.
The product-source system with symmetric
spatial coupling continues on every finite interval
by the energy estimate. Reconstruct $P$ and $Q$
from $P_\tau=V$ and $Q_\tau=e^{-P}X$,
with their initial values. The definitions of $W$
and $Y$ persist by the $W,Y$ equations and the initial
derivative identities. The usual
propagation of the momentum constraint gives
the periodic metric.

In the damped three-component subsystem,
the $X,Y$ spatial terms still cancel.
At order $s-1$ the remaining spatial forcing
of $W$ is bounded by $e^{-\tau}\|V-v_*\|_{H^s}$.
The nonlinear terms in these three equations
have two damped factors in the half-energy
estimate and a coefficient of size $2A_s\delta$.
Hence for some $\beta>0$ one has
$\|D(\tau_0+r)\|_{H^{s-1}}\le C e^{-\beta r}$.
This also follows by applying the scalar
damped differential inequality to
$\|D\|_{H^{s-1}}$ with input $e^{-\tau}\delta$.

For the difference energy at order $s$, write
$\Delta Z=Z-\widehat Z$. Every nonlinear term
in the half-energy estimate is bounded by
\[
 4A_s\delta\|\Delta D\|_{H^s}^2
  +3A_s(\|D\|_{H^s}+\|\widehat D\|_{H^s})
           \|\Delta Z\|_{H^s}\|\Delta D\|_{H^s}.
\]
The first term is absorbed by damping. Apply
Young's inequality to the second, retaining
damping and a coefficient bounded by
$C(\|D\|_{H^s}^2+\|\widehat D\|_{H^s}^2)$
on the difference energy.
This coefficient is time-integrable by
\eqref{eq:low-velocity-energy}. Consequently
$\sup_r\|\Delta Z\|_{H^s}\le C\mathcal I$.
The equation for $V_\tau$ now gives
$\|\Delta V_\tau\|_{H^{s-1}}
 \le C\mathcal I e^{-\beta r}$, reducing $\beta$
if necessary. The same integrable bound holds for
each solution's $V_\tau$ without $\mathcal I$.
In particular,
\[
 \|v-\widehat v\|_{H^{s-1}}
       +e^{\beta r}\|\Delta V(r)-(v-\widehat v)\|_{H^{s-1}}
 \le C\mathcal I .
\]
The limits of $P-rv$ and of $\lambda-rv^2$ follow
by integrating $V-v$ and
$\lambda_\tau-v^2=V^2-v^2+|D|^2$ respectively.
The differences of these integrands are
$O(\mathcal I e^{-\beta r})$ in $H^{s-1}$.

For the transverse limits, write
$P=P(\tau_0)+v_*r+B$, where
$\|B\|_{H^s}\le\delta r$.
The Sobolev algebra exponential series bounds
$\|e^{\pm B}\|_{H^s}\le C e^{A_s\delta r}$.
The difference of any two such exponentials
has the additional factor
$C(1+r)\mathcal I$, by the product formula
for the difference and
$\|P-\widehat P\|_{H^s}\le(1+Cr)\mathcal I$.
Thus $Q_\tau=e^{-P}X$ and its difference
are integrable in $H^s$ since $A_s\delta<v_*$.
The equation
$\Pi_\tau=e^{-2\tau}\partial_\theta(e^{2P}Q_\theta)$
and the uniform bound for $Q$ in $H^s$ give
integrable bounds for $\Pi_\tau$ and its difference in
$H^{s-2}$, using $A_s\delta<1-v_*$.
The difference estimate has at most an additional
polynomial factor in time. These estimates give the
two remaining limits and their Lipschitz bounds.

Finally $\lambda_\theta=2(V P_\theta+\Pi Q_\theta)$.
Subtract $2rv\,v_\theta$ and take the limit
in $H^{s-2}$ using the exponential remainders.
This gives \eqref{eq:tail-momentum}.
A periodic spatial primitive of the initial
momentum equation bounds the initial $\lambda$
norm by the initial orbit fields, velocities
and its value at a fixed angle.
\end{proof}

\begin{proposition}
\label{prop:from-observations-to-tail}
There are smooth inhomogeneous rank-two
backgrounds in the region of
Proposition~\ref{prop:low-velocity-difference}
for which the finite-snapshot data of
Theorem~\ref{th:fixed-rank-two} determine all five
asymptotic fields with Lipschitz dependence
in the stated norms. One can also use the
two-slice data of Proposition~\ref{prop:two-slice-open}
on an open set of such backgrounds. If neither
initial velocity is supplied, an explicit family
in the same low-velocity region admits recovery
from one full later orbit snapshot and a finite
set of early low-mode orbit coefficients.
\end{proposition}

\begin{proof}
Put $a=e^{-\tau_0}$. Start from
$P=v_*\tau+P_0$, $Q=$ constant.
Take a nearby nonvertical target geodesic
$\gamma$ and replace its scalar wave parameter
$v_*\tau+P_0$ by that parameter plus a small
nonconstant smooth even scalar wave.
The data at $a$ approach the vertical
homogeneous data in the required norms as both
perturbations vanish. Choose the geodesic and
wave so its tangent has nonzero $e_P$ component
on the range at $a$. This is the background
of Proposition~\ref{prop:two-slice-open}.
A small even perturbation of the initial
$Q_t$ transverse to the scalar wave direction
makes the target differential rank two on an open
arc where $P_\theta\ne0$. All orbit fields
and velocities at $a$ remain even, so the
momentum constraint still holds. The finite-time
inverse persists on a neighborhood. The tail
condition also persists by strict smallness.
Composition of the finite-time inverse with
Proposition~\ref{prop:low-velocity-difference}
gives the claim for fixed initial known fields.
Local variation of those fields is allowed as
well, by the upper Cauchy estimate and
the same fixed-background lower bound for
the missing velocity.

For the two-velocity case, at $a$ take
\[
 P=P_h(a)+\epsilon\cos\theta,\quad
 Q=\epsilon\sin\theta,\quad P_t=-v_*/a,\quad Q_t=0,
 \qquad P_h(t)=P_h(a)-v_*\log(t/a).
\]
These data satisfy the periodic momentum constraint.
They have rank two on open arcs and their
$\|Z(\tau_0)\|_{H^s}$ is $O(|\epsilon|)$.
At order one in $\epsilon$ the two displacement
perturbations are $p_1(t)\cos\theta$ and
$q_1(t)\sin\theta$. Their amplitudes start at one
with derivative zero at $a$.
At order two the orbit fields have Fourier modes
only up to order two. These claims follow by
differentiating the equations about the homogeneous
background once and twice.
Choose a sufficiently short rational delay
$d=\pi/L$ with $L>2$ so that $p_1q_1>0$ up to
$c=a+d$. The angle of the rotation in
Proposition~\ref{prop:both-ray-holonomy} has circle
mean $-2\int_a^c\overline{e^P Q_\theta}\,dt$.
In the sum over $L$ shifted rotations all nonzero
Fourier modes of order at most two cancel.
The zero mode of $e^PQ_\theta$ at second order
is $\epsilon^2 e^{P_h}p_1q_1/2$.
Consequently the closed-transport angle equals
\[
 -L\epsilon^2\int_a^c e^{P_h(t)}p_1(t)q_1(t)\,dt
   +O(\epsilon^3)
\]
uniformly on the circle. This is separated from
zero modulo $2\pi$ for small nonzero $\epsilon$.
The rational holonomy criterion and a suitably
early finite set of orbit coefficients give the
stable regular-time inverse for both velocities.
Proposition~\ref{prop:low-velocity-difference}
then gives the same asymptotic conclusion.
Alternatively the same initial curves satisfy the
independence condition of
Theorem~\ref{th:full-metric-gramian}, allowing a
short irrational delay when the terminal metric
gradient is also supplied.
\end{proof}

The estimate applies to the low-velocity region
defined at the recovered initial slice.
It does not assert control through spike
formation or a high-velocity transition.

\section{Stable frequency-assigned recovery in the nonlinear equations}
\label{sec:nonlinear-band}
The lower bounds at fixed snapshot times concern the information
record. Allowing each Fourier coefficient of a displacement field
to use a prescribed observation time changes that record. We now
give a nonlinear result valid in a regular-time neighborhood of
arbitrary smooth Gowdy data. The positive delay in the schedule is
independent of the spatial bandwidth.

Fix $a>0$, $s>1/2$ and a class of solutions on $[a,a+T]$ whose
$\|(P,Q)\|_{H^{s+1}}+\|(P_t,Q_t)\|_{H^s}$ are uniformly bounded.
Write $U=(P,Q)$ and write the coordinate field equations as
$U_{tt}-U_{\theta\theta}=\mathcal F(t,U,U_t,U_\theta)$,
where
\[
 \mathcal F=
 \begin{pmatrix}
 -P_t/t+e^{2P}(Q_t^2-Q_\theta^2)\\
 -Q_t/t-2(P_tQ_t-P_\theta Q_\theta)
 \end{pmatrix}.
\]
On this bounded class there is a Lipschitz constant $L$ such that
for the difference $Y=U-\widehat U$ and source difference $F$,
\[
 \|F(t)\|_{H^s}\le L E_s(Y(t)),\qquad
 E_s(Y)^2=\|Y_t\|_{H^s}^2+\|Y_\theta\|_{H^s}^2+\|Y\|_{H^s}^2.
\]
This follows from the algebra property of $H^s(\mathbb T)$ and
the Sobolev composition estimate for $e^{2P}$ on bounded sets.
For common initial $U(a)$, $E_s(Y(a))=\|Y_t(a)\|_{H^s}$.
The energy identity and Gronwall give
$E_s(Y(a+r))\le e^{(L+1)r}E_s(Y(a))$.
In particular let
\[
 \varepsilon_T = \frac{L}{L+1}\big(e^{(L+1)T}-1\big).
\]

For each recorded Fourier frequency $k$, use a delay $d_k\in(0,T]$
with $d_{-k}=d_k$ and set
\[
 w_k=\max\{|k|,1/d_k\},\qquad
 g_k=w_k S_k(d_k),\qquad
 S_k(r)=
 \begin{cases}
 \sin(|k|r)/|k|,& k\ne0,\\
 r,& k=0.
 \end{cases}
\]
We record $\widehat P(k,a+d_k)$, with the same initial
$P(a),Q(a),Q_t(a)$ for all compared solutions.
For comparing records, multiply each measured difference by $w_k$
and use the $H^s$ norm of the resulting Fourier sequence.
For a decoder, subtract the known free displacement
$\cos(|k|d_k)\widehat P(k,a)$ before multiplying.
The analogous record for both displacement fields
allows both initial velocities to be unknown.

\begin{theorem}[Nonlinear sampling estimate]
\label{th:nonlinear-band}
Suppose $\inf_k |g_k|\ge\gamma>\varepsilon_T$ on the
recorded frequencies. For two Gowdy solutions in the
bounded class, with common initial position and common
initial $Q_t$, write $f=P_t(a)-\widehat P_t(a)$,
and let $\mathcal M$ be the weighted record. With all
frequencies observed,
\[
 (\gamma-\varepsilon_T)\|f\|_{H^s}
 \le \|\mathcal M(U)-\mathcal M(\widehat U)\|_{H^s}
 \le (1+\varepsilon_T)\|f\|_{H^s}.
\]
With only $|k|\le K$ observed, the lower estimate becomes
\[
 (\gamma-\varepsilon_T)\|\Pi_K f\|_{H^s}
 \le \|\mathcal M_K(U)-\mathcal M_K(\widehat U)\|_{H^s}
       +\varepsilon_T\|(1-\Pi_K)f\|_{H^s}.
\]
The same estimates hold for both missing initial velocity
fields when both displacements are read.
\end{theorem}

\begin{proof}
Duhamel gives the observed $P$-difference on mode $k$
as $S_k(d_k)f_k+\int_0^{d_k} S_k(d_k-r)F_{P,k}(a+r)\,dr$.
For $0\le r\le d_k$, $w_k|S_k(d_k-r)|\le1$.
Indeed the two possible factors in $w_k$ use either
$|\sin x|\le1$ or $|S_k(r)|\le r$.
For the weighted sequence of integrals, use Minkowski after
extending each integral by zero to the common interval $[0,T]$.
The entire remainder is bounded in $H^s$ by
$\int_0^T\|F(a+r)\|_{H^s}\,dr\le\varepsilon_T\|f\|_{H^s}$.
This argument requires no diagonal structure for $F$ and admits
arbitrary coupling between modes.
The diagonal free term has gain between $\gamma$ and one.
The triangle inequality gives the two bounds and the version
with the omitted velocity tail. Vector-valued data use
the same argument. The periodic Gowdy momentum constraint
restricts the set of compared Cauchy data and does not
change the estimate.
\end{proof}

\paragraph{An explicit schedule with positive minimum delay.}
Set $N=\lceil 2\pi/T\rceil$ and $d_0=T/4$.
Read frequencies $|k|<N$ at $d_0$.
Their free phases have magnitudes below $\pi/2$, giving
$|g_k|\ge\sin(1)$, including $g_0=1$.
Group the remaining frequencies as
$N2^j\le |k|<N2^{j+1}$ for $j\ge0$.
For each group use the baseline $d_0$ and the option
$d_0+\pi/(4N2^j)$. For each mode keep the reading at
whichever of the two times gives the larger
$|\sin(|k|d)|$. Their phase separation is in
$[\pi/4,\pi/2)$, so the chosen magnitude is at least
$\gamma=\sin(\pi/8)$. These modes have $w_k=|k|$.
All chosen times are in $[a+T/4,a+3T/8]$.
There are at most $1+\max(0,1+\lfloor\log_2(K/N)\rfloor)$
distinct times through bandwidth $K$, and one saved
coefficient per frequency in addition to the initial
data. Since $1/d_k\le4/T$, the observation norm is
equivalent to $H^{s+1}$ on the assembled displacement
sequence, with constants independent of $K$.
At fixed $a$ and a fixed bounded set of nearby Cauchy data,
local existence gives a common regular-time interval and a common
bound, so $T$ can be reduced until $\varepsilon_T<\gamma$.
This includes neighborhoods of non-geodesic-image
solutions and the rank-two examples above.

\paragraph{Stable decoding and truncated data.}
Regard the forward map as $D v+R(v)$, with $D$ the
diagonal multipliers $g_k$ and $R$ Lipschitz with
constant $\varepsilon_T$. The subtraction of the known
free displacement has been included. Let $B$ be a closed
convex ball of trial velocities intersected with the
affine momentum constraint, and suppose the same forward
bound holds for every point of $B$.
For $|k|\le K$ take trial velocities in a fixed affine
band-limited slice, with their high-frequency part held fixed.
For both the infinite and finite-dimensional cases,
\[
 v_{n+1}=\operatorname{proj}_{B}
 \big[v_n+D^{-1}(\mathcal M_{\mathrm{data}}-\mathcal M(v_n))\big]
\]
is a contraction with ratio at most $\varepsilon_T/\gamma$.
In the affine finite-band case the multiplier inverse in
this formula acts only on the measured modes.
The projection is in the relevant Hilbert $H^s$ norm.
Consistent data recover the true velocity in $B$.
For additive error $\eta$ in the weighted record, the
resulting velocity error is at most
$\|\eta\|_{H^s}/(\gamma-\varepsilon_T)$ whenever the
true velocity lies in $B$. The truncated comparison
bound gives the additional $H^s$ tail term when the
true velocities are not in the same band-limited slice.
For an a priori $H^{s+r}$ bound, this tail is of order
$K^{-r}$ for $r>0$.

\paragraph{Timing errors.}
For known actual observation delays the proof applies
as written to those delays. For unknown timing errors
$\eta_k$ with both the nominal and actual times in
$[a,a+T]$, set $\eta_*=\sup_k w_k|\eta_k|$ using the
nominal weights. The change in the nominally weighted
record has norm bounded by
\[
 \eta_*\left(\|P_t(a)\|_{H^s}+\|P_\theta(a)\|_{H^s}
           +\int_0^T\|\mathcal F_P(a+r)\|_{H^s}\,dr\right).
\]
To see this, bound each temporal difference by
$|\eta_k|\sup_t|\widehat P_t(k,t)|$ and apply the
cosine propagator formula to $P_t$.
The bound is independent of the number of slots.
At high frequencies it corresponds to a clock error
on the order of $1/K$ for a uniform comparison.

\paragraph{Why the number of delayed slots grows.}
Consider a class containing the constant wave map and
suppose a uniform energy-normalized inverse estimate
holds with one later reading per mode and $m$ possible
delays in a fixed interval $[d_{\min},T]$, where
$d_{\min}>0$. Differentiate along a polarized single-mode
Cauchy family at the constant solution. The linear scalar
transfer in energy normalization is
$\sqrt{a/(a+d)}\{\sin(kd)+O(1/k)\}$,
uniformly on the fixed interval. For $k\ge 1/d_{\min}$
the prescribed weight is exactly $k$.
Choose a fixed integer $\ell$ large enough for the error
to be smaller than a quarter of the required lower gain
on multiples of $\ell$. Partition $[0,1)^m$ into
$R^m$ boxes for a fixed sufficiently large integer $R$.
Among $R^m+1$ multiples of the vector
$(\ell d_i/\pi)_{i=1}^m$ modulo one, two
belong to the same box. Their index difference gives
$k\in[\ell,\ell R^m]$ with $|\sin(kd_i)|\le\pi/R$
at every slot. This contradicts the required gain if
$\ell R^m\le K$. Thus the order $\log K$ is necessary
for uniform modewise energy conditioning in this class
when one reading per frequency is saved. The same
argument applies if the maximum of all available
readings per mode is used.

\paragraph{Spatial dispersion in a numerical decoder.}
A centered spatial stencil changes the free frequency
from $|k|$ to a known $\omega_h(k)$.
Even when $kh$ resolves the wave, a delayed slot incurs
phase error $d_k(\omega_h(k)-|k|)$.
For a stencil of order $p$ and all recorded modes
through $K$ with $Kh$ in the resolved range, a
sufficient condition to preserve the continuum phase
margin is $T K(Kh)^p\ll1$.
This differs from a requirement expressed only in
points per spatial wavelength.
For observations generated by a specified discretization
one can instead place the slots using its own
$\omega_h(k)$ in the two-phase rule.

A second option, when using a coarse inverse against
continuum-like data, corrects the modeled linear
transfer and retains the computed nonlinear remainder.
Write $\mathcal M_h(v)=C_h U(a)+D_h v+R_h(v)$,
where $C_h,D_h$ are the exact scalar damped-wave
transfers at the stencil frequencies, including the
observation weights. Replace the prediction by
\[
 \mathcal M_h^{\,\mathrm{corr}}(v)
    =\mathcal M_h(v)+(C-C_h)U(a)+(D-D_h)v .
\]
The correction uses the known initial displacement
and the current trial velocity.
If $\|R_h(v)-R_h(\widehat v)\|\le\epsilon_h\|v-\widehat v\|$
and the continuum diagonal $D$ has lower gain
$\gamma>\epsilon_h$, the correction makes the same
contraction proof apply to the inverse.
For a true continuum datum, the remaining error is
controlled by $R_h(v)-R(v)$ at the true initial
velocity, in addition to measurement noise.
Thus phase correction and nonlinear model bias are
separate. In the tests below the principal correction
removed the observed coarse-grid instability.

\input{archive_application.tex}
\section{A certified nonlinear basin on the archived grids}
\label{sec:certified-archive}
The earlier measured tangent near the boundary tape can be
replaced by a uniform finite-grid bound over a velocity ball.
Use the centered periodic stencil of the historical trace,
start from the depth-$18$ fields, and hold $P,Q,Q_\tau$ fixed
there. Write $V$ for the missing $P_\tau(18)$,
$V_*$ for the archived value, and
\[
 \mathcal F_{\mathrm{arch}}(V)
  =\frac{P(18)-P(16;V)}2,\qquad
 Y_{\mathrm{arch}}=\frac{P_{\mathrm{archive}}(18)-
                              P_{\mathrm{archive}}(16)}2 .
\]
The archived velocity is used to locate a proof ball.
It is not needed to run the reconstruction.

\begin{theorem}[Archive certificate]
\label{th:archive-certificate}
For each $N=1024,2048,4096$, the exact centered-stencil
ordinary differential equation initialized by the stored
floating-point arrays is well defined throughout the
window for all $\|V-V_*\|_\infty\le0.06$.
On that entire ball,
\[
 \|D\mathcal F_{\mathrm{arch}}(V)-I\|_\infty < 0.10725 .
\]
The recorded $Y_{\mathrm{arch}}$ lies in the image of the
ball, and there is a unique compatible velocity in it.
The iteration
\[
 V^{(0)}=Y_{\mathrm{arch}},\qquad
 V^{(m+1)}=V^{(m)}+Y_{\mathrm{arch}}
                         -\mathcal F_{\mathrm{arch}}(V^{(m)})
\]
starts in the ball and converges to that velocity.
Additional noise in the normalized record of supremum
norm at most $0.05095$ still gives a unique solution in
the same ball, with velocity error at most $1.121$
times the record noise. The certificate uses the exact
semidiscrete equations for each stated $N$.
\end{theorem}

\begin{proof}
Let $x=18-\tau\in[0,2]$ and denote spatial difference
operators by $D_1,D_2$. At each site set
\[
 B=e^P Q_\tau,\qquad C=e^{P-\tau}D_1Q,\qquad
 Z_i=e^{P_i-\tau}\frac{Q_{i+1}-Q_i}{\Delta\theta},
 \qquad \xi=e^{-\tau}/\Delta\theta .
\]
The exact reduced equations give
\begin{align*}
 V_x&=C^2-B^2-e^{-2\tau}D_2P,\\
 B_x&=VB-e^{P-2\tau}(D_2Q+2D_1P D_1Q),\\
 C_x&=(1-V)C-e^{P-\tau}D_1Q_\tau,\\
 (Z_i)_x&=(1-V_i)Z_i-\xi
          (e^{P_i-P_{i+1}}B_{i+1}-B_i).
\end{align*}
Bootstrap $0.31\le V\le1.425$,
$|B|\le0.035$, $|C|\le0.276$ and $|Z|\le0.965$.
The adjacent $P$ increments are bounded by
$0.18+2(1.425-0.31)=2.41=:J$.
With $E=e^J$, the spatial term in the $V$ equation is
bounded by $2\xi^2 J$. The absolute forcing bounds
for $B,C,Z$ after their displayed local terms are
respectively
\[
 \xi\big((1+E)|Z|+2J|C|\big),\qquad
 \xi E|B|,\qquad \xi(1+E)|B|.
\]
These follow directly from the two neighboring edge
differences. Integrating the scalar growth inequalities
bounds $V,B,Z$ strictly inside their proposed limits.

For $C$, compare the pair $(u,C)=(1-V,C)$ to the
sitewise velocity-dominated system
$u_x=-C^2$, $C_x=uC$ with the same initial values.
This system conserves $u^2+C^2=\sigma^2$ and admits the
explicit transfer
\[
 C_{\mathrm{ref}}(x)
 =\frac{C(0)}{\cosh(\sigma x)-(u(0)/\sigma)\sinh(\sigma x)}.
\]
When $u$ starts nonpositive the amplitude decreases.
Otherwise its maximum occurs at the end of the window
or when $u=0$, and the growth estimate
$|C_{\mathrm{ref}}(x)|\le |C(0)|e^{x\max(u(0),0)}$
holds throughout. These bounds provide the interval
maximum for each possible initial $V$ in the ball.
The comparison Jacobian in $(u,C)$ has maximum row
norm bounded by
$\max\{2|C|,|C|+|u|\}$ on the segment between the
two solutions. If $L$ and $f$ bound that norm and
the forcing in the two-component comparison, its
maximum error at backward time $x$ is
$f(e^{Lx}-1)/L$. The resulting bound on $C$ strictly
improves the proposed cap. Continuation closes the
bootstrap for every velocity in the ball.

For the tangent derivative, use the scaled perturbations
$(p,v,y,z)=(\delta P,\delta P_\tau,e^P\delta Q,
e^P\delta Q_\tau)$, with $p=y=z=0$ initially.
The component supremum norms satisfy a comparison
inequality with a constant nonnegative matrix $M$.
Writing $B_*,C_*,V_{\max}$ for the established bounds
and $\xi_*=\max_{[0,2]}\xi$, the rows of one such $M$ are
\[
 \begin{pmatrix}
 0&1&0&0\\
 2\max(B_*^2,C_*^2)+4\xi_*^2&0&2\xi_*C_*E&2B_*\\
 0&0&0&1\\
 2\xi_*C_*&2B_*&[2+2E(1+J)]\xi_*^2&V_{\max}
 \end{pmatrix}.
\]
For example, the $\xi_*^2$ terms come from the centered
second difference in $P$ and the corresponding weighted
difference in $Q$. The term $-V y$ in the backward
equation can be dropped for this upper comparison.
Since $p_x=-v$, bounding the feedback into $v$ gives
\[
 \|D\mathcal F_{\mathrm{arch}}-I\|_\infty
 \le \frac12 (e^{2M}e_v)_p-1 ,
 \qquad e_v=(0,1,0,0)^{\mathsf T}.
\]
The subtraction here removes the free contribution
of the initial velocity. In detail the difference from
the free position is the time integral of the $v$
feedback with kernel $2-x$, bounded by the same
integral in the positive comparison system.

The normalized position transfer of the sitewise
comparison at $V_*$ is
\[
 Y_{\mathrm{ref}}=
 1+\tfrac12\log\!\left(\cosh(2\sigma)
       -(u(0)/\sigma)\sinh(2\sigma)\right).
\]
Integrating the $u$ error above bounds
$\|\mathcal F_{\mathrm{arch}}(V_*)-Y_{\mathrm{ref}}\|_\infty$.
The measured $Y_{\mathrm{arch}}-Y_{\mathrm{ref}}$
is enclosed separately using the stored arrays.
The interval script checks the bootstrap,
evaluation of $Y_{\mathrm{ref}}$, and the matrix
exponential by positive terms through order $24$ with
a geometric tail bound. It uses interval evaluations
of the initial floating-point arrays and records
their SHA-256 hashes. Its results are
\[
\begin{array}{rccc}
 N&\|D\mathcal F_{\mathrm{arch}}-I\|_\infty
   &\|\mathcal F_{\mathrm{arch}}(V_*)-Y_{\mathrm{arch}}\|_\infty
   &\|Y_{\mathrm{arch}}-V_*\|_\infty\\ \hline
 1024&<0.10310&<0.002610&<0.04681\\
 2048&<0.10402&<0.002610&<0.04708\\
 4096&<0.10725&<0.002610&<0.04708 .
\end{array}
\]

The exact forward-based iteration is a contraction
on a complete ball with uniform factor $0.10725$.
For its self-map condition, the distance of the image
of $V_*$ from $V_*$ is the record mismatch.
The mismatch plus record noise must be at most
$(1-0.10725)0.06$. The interval bounds give the
stated noise budget and verify that the projected
secant itself lies inside the ball without noise.
Two compatible solutions in the ball obey the lower
Lipschitz estimate with factor $1-0.10725$.
This completes the certificate.
\end{proof}

The admissible-noise bound is a bound on the existence
and uniqueness of the nearby decoded velocity.
An arbitrary noisy secant need not serve as an initial
iterate in the same ball; the original secant or a
previous iterate in the certified ball can be used.
The argument treats the exact finite-dimensional
ordinary differential equations. In particular it
does not impose transfer of these constants to the
continuum or to the fourth-order stencil.


\paragraph{All four recorded fields.}
The same interval comparison can vary the stored
$P,Q$ and momentum $M=e^{2P}Q_\tau$ as well as the
normalized record. About the archived initial fields,
allow
\[
 \|\Delta P\|_\infty\le0.002,\quad
 \|e^{P_*}\Delta Q\|_\infty\le0.003,\quad
 \|e^{-P_*}\Delta M\|_\infty\le0.00015 .
\]
Along a straight segment in these three measured
variables, the scaled tangent of $Q_\tau$ begins with
$e^{-P}\Delta M-2B\Delta P$ and the scale for
$\Delta Q$ changes by at most $e^{0.002}$.
The same positive comparison system therefore
bounds the change in the normalized record at fixed
velocity by $0.000347$ over the three grids.
With these measured-field uncertainties, the uniform
derivative defect remains below $0.10866$ and
a further normalized secant error of $0.05052$
fits inside the certificate. The stated norms
make explicit how precision in the orbit momentum
depends on the marked $P$ scale.

\section{Estimating the limiting metric modulus from earlier data}
\label{sec:metric-modulus}
The polarized Einstein--scalar benchmark distinguishes the
additive metric datum from the derivative and shape fields.
Its production initial profile is a finite Fourier sum.
This permits an independent continuum computation of the
limiting modulus, and clarifies a more general stability
question for finite observations. The logarithmic term in
the field intercept is the established Kasner scattering
correction \cite[Theorem~1.3 and Section~1.4]{li}.
The result below concerns its contribution to the nonlinear
pointwise metric modulus after inversion of two field
snapshots.

Use the positive-quadrature convention with
$\tau=-\log t$ and scalar $\phi=\sqrt2\varphi$, where
$\varphi$ is canonically normalized at $8\pi G=1$.
The orbit field $P$ and scalar $\phi$ each satisfy
$u_{tt}+t^{-1}u_t-u_{\theta\theta}=0$.
The metric coefficient transverse to the two
Killing directions is $t^{-1/2}e^{-\lambda/2}$.
With $U=(P,\phi)$ one has
$\lambda_\tau=|U_\tau|^2+e^{-2\tau}|U_\theta|^2$.
For a fixed angular origin define
$v=\lim_{\tau\to\infty}U_\tau$,
$U_\infty=\lim_{\tau\to\infty}(U-\tau v)$,
and $c_\lambda=\lim_{\tau\to\infty}(\lambda(\tau,0)-\tau|v(0)|^2)$.
The origin can be replaced by any fixed angle.

At $t=a$, with $\tau_a=-\log a$ and one measured metric anchor,
\begin{align}
 c_\lambda
 &=\lambda(\tau_a,0)-\tau_a|v(0)|^2
   +\int_{\tau_a}^\infty
     \left(|U_\tau|^2-|v|^2+e^{-2\tau}|U_\theta|^2\right)(\tau,0)
     \,d\tau \nonumber\\
 &=\lambda(\tau_a,0)+\tau_a|v(0)|^2
     +2v(0)\cdot\big(U_\infty(0)-U(\tau_a,0)\big)
     +\int_{\tau_a}^\infty
        \left(|U_\tau-v|^2+e^{-2\tau}|U_\theta|^2\right)(\tau,0)
        \,d\tau .
 \label{eq:modulus-identity}
\end{align}
The second equality follows by integrating the part
linear in $U_\tau-v$. It displays the dependence on
the intercept as well as the slope. Both formulas
use one finite metric anchor, whose additive change
adds the same constant to $c_\lambda$.

\subsection{A sharp logarithmic threshold for the missing modulus}
For $\alpha\in\mathbb R$, let $\mathcal H^2_\alpha$
denote the space with Fourier coefficient weight
$\langle k\rangle^2[\log(2+|k|)]^\alpha$ in $\ell^2$.
Consider two full orbit-and-scalar snapshots at
$t=a,b$, with rational $b-a$ in units of $\pi$
and no zero modal transfer at that pair.
The excluded reference times form a locally finite
set at each such gap, by the scalar phase estimate.

\begin{theorem}[Sharp threshold for the pointwise modulus]
\label{th:metric-modulus}
With one metric anchor at the fixed angle and time $a$,
the limiting metric modulus in the polarized
Einstein--scalar system is locally Lipschitz in the
two snapshots with the $\mathcal H^2_\alpha$ norm if
$\alpha>3/2$. In particular it is locally Lipschitz
in $H^r$ for every $r>2$.
Near a homogeneous datum with nonzero limiting
velocity, a local Lipschitz estimate fails on smooth
periodic, constraint-compatible data for
$\alpha\le3/2$, and hence in the unweighted
Sobolev space $H^2$ or lower.
\end{theorem}

\begin{proof}
For either field, write the nonzero modes as
$A_kJ_0(|k|t)+B_kY_0(|k|t)$.
At the stated observation pair the inverse Bessel
matrix gives an upper bound
\[
 |A_k|+|B_k|\le C\langle k\rangle^{3/2}
                 \big(|\widehat U(a,k)|+|\widehat U(b,k)|\big).
\]
Indeed the Hankel expansion for $\beta=1/4$
gives order $k^{-1}$ for the energy-normalized
transfer on the resonant subsequence and a uniform
lower bound away from that subsequence. The finite
nonzero transfers and the scalar zero mode pose
no further difficulty. Thus the snapshots bound
the coefficient norm with weight
$\langle k\rangle^{1/2}[\log(2+|k|)]^\alpha$.
The limiting mode formulas are
\[
 v_k=-\frac{2B_k}{\pi},\qquad
 (U_\infty)_k=A_k+
      \frac{2B_k}{\pi}\big(\log(|k|/2)+\gamma_{\mathrm E}\big).
\]
Their absolute Fourier sums are bounded for
$\alpha>3/2$, by Cauchy--Schwarz.

We show that the integrals in the second line of
\eqref{eq:modulus-identity} are bounded quadratic
forms in the same data norm. At fixed $\theta$,
in the measure $dt/t$, consider
$G=U_\tau-v$ and $E=tU_\theta$.
For $z=|k|t\le1$, Bessel estimates give
\[
 |G_k|\le C(|A_k|+|B_k|)z^2(1+|\log z|),\qquad
 |E_k|\le C(|A_k|+|B_k|)z(1+|\log z|).
\]
The $L^2(dt/t)$ norm of each bound is at most a
constant times the coefficient magnitude,
uniformly in $k$. Their sum is finite in the stated
coefficient norm.

For $z\ge1$, the leading terms of both
$U_\tau$ and $E$ are sums of
$\sqrt{|k|t}\,c_k^\pm e^{\pm i|k|t}$
with $|c_k^\pm|\le C(|A_k|+|B_k|)$.
After removing the common $\sqrt t$, untruncated
sums are bounded in $L^2(0,a;dt)$ by Fourier
orthogonality over periods of length $2\pi$.
The missing intervals $t<1/|k|$ are controlled by
the matrix $1/\max(k,\ell)$ on $\ell^2$ for positive
integers, applied to coefficients $\sqrt{k}|c_k^\pm|$.
Its boundedness follows from Schur's test with
weight $k^{-1/2}$. The same matrix controls
the large-argument $O(z^{-1/2})$ remainders
because $\int_{\max(1/k,1/\ell)}^a t^{-2}\,dt
\le\min(k,\ell)$. Finally, the subtraction of $v_k$
on $t\ge1/|k|$ costs at most
$C|B_k|\sqrt{\log(2+|k|)}$ per mode in $L^2(dt/t)$.
Its sum is bounded already for $\alpha>1$.
These estimates give uniform pointwise-in-$\theta$
bounds on $G,E$ in $L^2(dt/t)$.
By approximation from finite Fourier data,
the limiting formula and its local Lipschitz bound
follow for $\alpha>3/2$.

For sharpness, linearize \eqref{eq:modulus-identity}
at a homogeneous field with nonzero velocity in at
least one component, using an even perturbation in
that component and fixed metric anchor.
The quadratic integral contributes zero at first order.
The derivative of the first line reduces to
$2v\cdot(\delta U_\infty-\delta U(a))$ at the chosen
angle. Take $\delta U(a)=0$.
On the resonant multiples of the rational
denominator, Cramer's rule gives
\[
 B_k=J_0(ka)\widehat{\delta U}(b,k)/
       [J_0(ka)Y_0(kb)-Y_0(ka)J_0(kb)] .
\]
On a positive-density subsequence the asymptotic cosine
of $J_0(ka)$ has magnitude bounded below;
this follows from periodicity or equidistribution
of the rotation. On these modes the limiting
intercept has multiplier of order $k^{3/2}\log k$.
Translation of the even perturbations to the
measurement angle preserves the argument.
They satisfy the periodic momentum constraint since
their spatial derivative is odd about that angle
and their initial position was fixed and homogeneous.
A bounded derivative acting on the snapshot
$\mathcal H^2_\alpha$ norm would require convergence
of a sum comparable to
$\sum k^{-1}(\log k)^{2-2\alpha}$ on the indicated
subsequence. This fails for $\alpha\le3/2$.
Taking finite real Fourier perturbations and
differentiating any proposed local Lipschitz estimate
shows its failure on smooth admissible data.
\end{proof}

\subsection{Averaging the metric datum changes its critical norm}
\label{sec:modulus-aperture}

Let $\chi\in L^\infty(\mathbb T)$ be a nonnegative
normalized spatial weight of positive measure, and
write $\langle\chi,f\rangle=\int_{\mathbb T}\chi f$.
Supply the averaged anchor $\langle\chi,\lambda(\tau_a)\rangle$.
The corresponding absolute metric datum is the limit
$\langle\chi,\lambda_\infty\rangle$.
Averaging \eqref{eq:modulus-identity} determines it
from the same field snapshots.

\begin{proposition}
\label{prop:averaged-modulus}
For the fixed rational observation gap of
Theorem~\ref{th:metric-modulus}, the averaged modulus
is locally Lipschitz in the unweighted $H^2$ snapshot
norm. For every $r<2$ it fails to be continuous
near homogeneous data in the $H^r$ snapshot norm.
If the two fields instead use stable frequency-assigned
displacement measurements in the usual wave-energy
normalization, the corresponding sharp averaged
threshold is $H^1$. In that design the sharp
pointwise modulus threshold is
$\mathcal H^1_\alpha$ with $\alpha>3/2$,
where the Fourier weight is
$\langle k\rangle[\log(2+|k|)]^\alpha$.
\end{proposition}

\begin{proof}
For a nonzero Bessel mode, set $z=|k|t$ and write
$E_k=t\,\widehat U_\theta$. The large-argument bounds
give
\[
 \int_{\tau_a}^\infty
      (|(U_\tau-v)_k|^2+|E_k|^2)\,d\tau
 \le C\langle k\rangle(|A_k|^2+|B_k|^2).
\]
One may integrate the leading $O(z)$ energy
over $1\le z\le |k|a$, the constant subtraction
over the same interval, and the integrable
small-argument remainders over $z<1$.
Parseval therefore controls the time integral in
the averaged identity at any bounded $\chi$
by the coefficient norm with weight
$\langle k\rangle^{1/2}$.
The products involving $v,U_\infty$ are bounded
in the same coefficient norm by the spatial
$L^2$ Cauchy--Schwarz inequality, since
$(\log|k|)^2=O(\langle k\rangle)$.
For a fixed rational two-snapshot gap this
coefficient norm is bounded by the snapshot $H^2$
norm. The averaging functional is quadratic plus
the supplied anchor, so the same estimates give
local Lipschitz continuity.

For sharpness, choose a single even perturbation
of a homogeneous field with zero datum at $a$.
Let its only later datum have amplitude $k^{-2}$
on a resonant multiple of the denominator.
Its Bessel coefficients have squared norm
of order $k^{-1}$, since the Bessel determinant
is of order $k^{-2}$ and
$J_0(ka)^2+Y_0(ka)^2\asymp k^{-1}$.
In the mean, the principal large-argument
energy has integrand
$(2/\pi)z(|A_k|^2+|B_k|^2)$ up to the fixed
real Fourier normalization.
This follows by adding the two squared
quadrature phases of the temporal and angular
terms. Its time integral is of order
$k(|A_k|^2+|B_k|^2)$, whereas the subtracted
limiting energy and the lower-order terms
contribute at most
$O(\log k)(|A_k|^2+|B_k|^2)$.
For $\chi$ the same leading contribution
survives: $\cos^2(k\theta)$ and
$\sin^2(k\theta)$ each have averaged value
$1/2+o(1)$ by the Riemann--Lebesgue lemma.
The linear term about the homogeneous
field is of size
$O(k^{-1/2}\log k)$ even without angular
averaging. Thus the modulus changes by a
quantity bounded away from zero while the
$H^r$ data tend to zero for $r<2$.

Under a stable frequency-assigned scheme, the
one-derivative resonant loss in the fixed-pair
inverse is absent. The Bessel coefficient
weight $\langle k\rangle^{1/2}$ is then
bounded by the assigned displacements in $H^1$.
The same proof with amplitude $k^{-1}$ for the
later data gives failure below $H^1$.
The pointwise proof of
Theorem~\ref{th:metric-modulus} uses the same
coefficients with an extra logarithmic weight, so
its upper threshold shifts down by one derivative.
For its lower bound the normalized scalar gains
are bounded above as well as below. Along a
positive-density set of modes with
$|J_0(ka)|\gtrsim k^{-1/2}$ the intercept
multiplier has order $k^{1/2}\log k$.
The same logarithmic sum gives the claimed
sharp endpoint.
\end{proof}

For a fixed smooth aperture profile $\chi_1$ of mass
one, its periodized width-$\delta$ rescaling
$\chi_\delta$ has rapidly decreasing Fourier
weights past frequency $1/\delta$.
At the homogeneous background of
Theorem~\ref{th:metric-modulus}, the $H^2$ norm of the
linearized fixed-pair averaged-modulus functional
satisfies
\[
 \|D\langle\chi_\delta,\lambda_\infty\rangle\|
      \asymp |\log\delta|^{3/2}
      \quad\text{as }\delta\downarrow0 .
\]
Indeed, on the indicated resonant subsequence the
squared norm is comparable to the partial sum
$\sum_{k\lesssim1/\delta}(\log k)^2/k$.
The aperture multiplier stays bounded below on a
fixed inner part of that range since its Fourier
transform equals one at zero. Its rapid decay
controls the remaining modes. This describes the
linearized cost of approaching the pointwise
metric convention with a narrowing spatial aperture.

\paragraph{An exact expression for the circle mean.}
For normalized integration over the circle write
\[
 E=\overline{|U_\tau|^2+e^{-2\tau}|U_\theta|^2},
 \qquad G=\overline{U\cdot U_\tau},
 \qquad \bar\lambda=\overline{\lambda}.
\]
In the polarized equations, periodic integration by parts yields
\[
 E'=-2e^{-2\tau}\overline{|U_\theta|^2},\qquad
 G'=\overline{|U_\tau|^2-e^{-2\tau}|U_\theta|^2}.
\]
Since $\bar\lambda'=E$, the combination
$\bar\lambda+E-G$ is constant. Evaluating it at
$\tau_a$ and in the asymptotic expansion gives the
quadrature-free identity
\begin{equation}
 \overline{\lambda_\infty}
 =\bar\lambda(\tau_a)+E(\tau_a)-G(\tau_a)
       -\overline{|v|^2}+\overline{U_\infty\cdot v}.
 \label{eq:mean-modulus-virial}
\end{equation}
For full vacuum wave maps, use
$G=\overline{P P_\tau}$ and the same wave-map energy
$E=\overline{|U_\tau|_{\mathbb H^2}^2+
                 e^{-2\tau}|U_\theta|_{\mathbb H^2}^2}$.
The exact derivative is instead
\[
 (\bar\lambda+E-G)'=
 \overline{(1-P)e^{2P}(Q_\tau^2-e^{-2\tau}Q_\theta^2)} .
\]
This identity holds at regular times. In an AVTD class
where the stated limits of $E,G$ and $\bar\lambda$ exist
and the right-hand side is integrable, it gives the
corresponding mean limit by adding its time integral.
The polarized identity requires no such measurement of
the transverse source.

\section{Numerical checks and the production metric modulus}
\label{sec:continued-numerics}
These computations check specific formulas and finite-grid
behavior. The continuum estimates above follow from the
proved bounds.

\begin{figure}[htbp]
\centering
\begin{tikzpicture}
\begin{axis}[width=.47\textwidth,height=.32\textwidth,
  xlabel={Fraction of the observation window},
  ylabel={Relative orbit displacement},
  ylabel style={font=\scriptsize},xlabel style={font=\small},
  tick label style={font=\scriptsize},
  ymin=-.27,ymax=.28,xmin=0,xmax=1,
  legend style={font=\scriptsize,at={(.5,.03)},anchor=south,
                legend columns=2,draw=none,fill=none}]
\addplot[teal!75!black,thick] table[x=fraction,y=branch_a,col sep=comma]
 {figures/collision_path.csv};
\addlegendentry{First history}
\addplot[orange!85!black,thick,dashed] table[x=fraction,y=branch_b,col sep=comma]
 {figures/collision_path.csv};
\addlegendentry{Second history}
\addplot[gray,densely dotted,forget plot] coordinates {(.2666666667,-.27) (.2666666667,.28)};
\end{axis}
\end{tikzpicture}\hfill
\begin{tikzpicture}
\begin{loglogaxis}[width=.47\textwidth,height=.32\textwidth,
  xlabel={Fourier frequency $k$},ylabel={Normalized orbit gain},
  xlabel style={font=\small},ylabel style={font=\small},
  tick label style={font=\scriptsize},
  legend style={font=\scriptsize,at={(.03,.03)},anchor=south west,draw=none,fill=none}]
\addplot[orange!85!black,thick,mark=square*] table[x=mode,y=nodal,col sep=comma]
 {figures/null_gains.csv};
\addlegendentry{Nodal slice}
\addplot[teal!75!black,thick,mark=*] table[x=mode,y=placed,col sep=comma]
 {figures/null_gains.csv};
\addlegendentry{Off-node slice}
\addplot[gray,densely dotted,forget plot] table[x=mode,y=principal,col sep=comma]
 {figures/null_gains.csv};
\end{loglogaxis}
\end{tikzpicture}
\caption{An exact rank-two collision on the marked torus.
The left panel shows $\widetilde P$ at one angle after
subtracting the common homogeneous background.
The curves match at both ends. The dotted line marks
the replacement terminal slice.
The right panel compares the linearized orbit response
on carriers that are in free resonance at both locations.
The response decays at the original nodal slice and
approaches the nonzero transport prediction (dotted)
at the replacement slice.}
\label{fig:null-numerics}
\end{figure}

For Fig.~\ref{fig:null-numerics}, the coordinate tangent
equations were integrated independently of the null-ray
transfer formula. The two ray angles agreed with
direct line integration within $9\times10^{-16}$.
The off-node normalized gain tends to the predicted
$0.08559$ on the tested sequence, whereas the nodal
gain falls from $0.03133$ at $k=64$ to $0.002041$
at $k=1024$. At bandwidth $96$, two grids gave the
minimum full-orbit singular gain $0.04090$ at the
replacement slice and gains below $5\times10^{-13}$
at the original nodal slice. A separate smaller-amplitude
test verified quadratic variation of the $P$-operator
and linear variation of the $Q$-operator as predicted
in Theorem~\ref{th:exact-rank-recovery}.

For nonlinear inversion at fixed times, an independent
higher-grid Fourier evolution supplied a perturbed
rank-two record. At bandwidth $96$, the two inverse
resolutions agreed to the shown precision.
Adding a normalized record error $3\times10^{-5}$ in
one resonant $P$ mode gave a recovered velocity
$H^1$ error $3.62\times10^{-4}$ from the full orbit
fields. Reading $P$ alone gave an $H^1$ velocity
error $0.06022$ despite a final position residual
below $6\times10^{-11}$. With exact data the full
orbit reconstruction reached a velocity error of
about $5.3\times10^{-10}$.
Reading $P$ and the terminal metric gradient in the
same noisy run gave $5.21\times10^{-6}$ velocity
error in the tested direction. These comparisons use
the stated Fourier-weighted norms and a single, specified
noise direction.

\begin{figure}[htbp]
\centering
\begin{tikzpicture}
\begin{loglogaxis}[width=.68\textwidth,height=.30\textwidth,
  xlabel={Fourier frequency $k$},ylabel={Minimum normalized gain},
  ymin=2e-5,ymax=.29,tick label style={font=\scriptsize},
  xtick={153,467,624,1575,2352},xticklabels={153,467,624,1575,2352},
  legend style={font=\scriptsize,at={(.03,.03)},anchor=south west,draw=none}]
\addplot[orange!85!black,thick,mark=square*] table[x=mode,y=orbit,col sep=comma]
 {figures/full_metric_gains.csv};
\addlegendentry{Orbit fields}
\addplot[teal!75!black,thick,mark=*] table[x=mode,y=metric,col sep=comma]
 {figures/full_metric_gains.csv};
\addlegendentry{Orbit fields and metric gradient}
\end{loglogaxis}
\end{tikzpicture}
\caption{Both velocities are missing in the circular-loop
example at an irrational delay. The plotted modes sample
near cancellations of the orbit-only transfer.
The metric gradient at the same terminal slice
raises the normalized gain.}
\label{fig:irrational-metric}
\end{figure}

Figure~\ref{fig:irrational-metric} compares smallest
singular gains for both velocities at sampled modes on
a circular-loop background at $d=0.02$.
At $k=1575$ the orbit-only gain is about
$4.23\times10^{-5}$ and the gain with the metric
gradient is about $0.118$. At the longer
illustrative gap $d=\pi\sqrt2/24$, the corresponding
gains at $k=3356$ are about $1.4\times10^{-5}$
and $0.678$. These are modal numerical comparisons.
The all-frequency short-delay claim in
Proposition~\ref{prop:irrational-metric} follows
from the analytic reverse estimate.

For the family of Proposition~\ref{prop:from-observations-to-tail},
the computed holonomy angle divided by
$\epsilon^2$ was $-2.94478497$ at
$\epsilon=2.5\times10^{-4}$, against the
quadratic prediction $-2.94478495$.
Two spatial grids gave depth-28 estimates of
the five asymptotic quantities agreeing to
$8\times10^{-14}$ in the stated norms.
The limiting momentum identity had residual
$1.24\times10^{-14}$ in $H^1$ on the finer grid.
These finite-depth checks illustrate the two
steps of the result and the asymptotic constraint.

The frequency-assigned experiment used a separate
rank-two pair for which the two original orbit snapshots
agreed within $1.3\times10^{-12}$ in direct forward
evolution. The delayed sampling resolved their
initial-velocity difference and had a normalized gain
of about $0.891$ on that pair. At bandwidth $288$ four
delayed slots gave a minimum free gain $0.529$,
compared with $2.46\times10^{-4}$ using the baseline
slot alone. The nonlinear reconstruction used an
independent fourth-order spatial stencil. At that
bandwidth it failed to converge on the coarsest
inverse grid and the recovered velocity improved
on the two finer grids. The mismatch is consistent
with stencil dispersion at inadequate spatial
resolution. Unknown timing offsets within the tested
$1/K$ scale introduced visible bias, which was reduced
when the true timestamps were used.
Substituting the exact continuum damped-wave transfer
for the stencil's free transfer removed the principal
dispersion discrepancy. At bandwidth $288$ on the
coarse grid the corrected inverse gave an $H^1$
velocity error $3.60\times10^{-9}$ in two updates
on the small-amplitude test. A more strongly coupled
profile retained a nonzero nonlinear model discrepancy
after the same correction. This separates finite-grid
phase error from nonlinear model error.

\begin{figure}[htbp]
\centering
\begin{tikzpicture}
\begin{semilogyaxis}[width=.47\textwidth,height=.32\textwidth,
  xlabel={Iteration from the secant},ylabel={Normalized record residual},
  xlabel style={font=\small},ylabel style={font=\small},
  tick label style={font=\scriptsize},
  legend style={font=\scriptsize,at={(.97,.98)},anchor=north east,draw=none,fill=none}]
\addplot[teal!75!black,thick,mark=*] table[x=iteration,y=residual,col sep=comma]
 {figures/archive_iteration.csv};
\addlegendentry{Archived grid $N=4096$}
\addplot[gray,densely dashed] table[x=iteration,y=certified_upper,col sep=comma]
 {figures/archive_iteration.csv};
\addlegendentry{Certified contraction bound}
\end{semilogyaxis}
\end{tikzpicture}\hfill
\begin{tikzpicture}
\begin{loglogaxis}[width=.47\textwidth,height=.32\textwidth,
  xlabel={Spatial grid size},ylabel={Absolute error in $c_\lambda$},
  xtick={256,512,1024,2048},xticklabels={256,512,1024,2048},
  xlabel style={font=\small},ylabel style={font=\small},
  tick label style={font=\scriptsize},
  legend style={font=\scriptsize,at={(.98,.54)},anchor=north east,draw=none,fill=none}]
\addplot[orange!85!black,thick,mark=square*] table[x=grid,y=zero_tail_depth8,col sep=comma]
 {figures/modulus_convergence.csv};
\addlegendentry{Earlier anchor, zero tail}
\addplot[teal!75!black,thick,mark=*] table[x=grid,y=bessel_tail_depth8,col sep=comma]
 {figures/modulus_convergence.csv};
\addlegendentry{Earlier anchor, Bessel tail}
\end{loglogaxis}
\end{tikzpicture}
\caption{Left, the archived reconstruction error in the
normalized record during the secant-started iteration,
compared to the factor from Theorem~\ref{th:archive-certificate}.
Right, the depth-eight smooth production metric-modulus
estimate with the finite-depth anchor at $\tau=6$ and
both field snapshots earlier than the discarded endpoint
at depth eight. The error is against the independent
continuum reference.}
\label{fig:archive-modulus}
\end{figure}

On all three archived centered grids the secant-started
replay took nine updates to reach a normalized residual
below $6\times10^{-13}$. At $N=4096$ it returned to the
held-out velocity within $4.1\times10^{-12}$ in
supremum norm. The numerical replay used the same
stencil as the archived tape. The separate interval
certificate guarantees a contraction in a whole velocity
ball for the exact semidiscrete evolution. A second
interval calculation permits simultaneous variation
of the recorded $P,Q$, weighted momentum and
normalized secant, at the stated radii.

The earlier modulus checklist proposed using the
depth-16 versus depth-24 difference as a spatial
convergence test. An independent Fourier--Bessel
calculation from the production initial data gives
$c_\lambda\simeq0.1180324739805121$ in the historical
anchor normalization, with an analytic tail bound
below $2.1\times10^{-28}$ after $\tau=34$.
At depth eight, the pre-endpoint, zero-tail estimate
has error close to $3.79\times10^{-4}$ on the fine grids.
Using the exact Bessel tail with the same earlier
anchor and two numerical field snapshots gives,
over $N=256,512,1024,2048$, errors close to
$8.93\times10^{-6},\,2.23\times10^{-6},\,
5.58\times10^{-7},\,1.40\times10^{-7}$.
The analytic semidiscrete reference reproduces those
spatial errors. The pre-endpoint depth-16 and
depth-24 zero-tail estimates differ by about
$2.3$ to $2.8\times10^{-10}$ over the grids,
without second-order decrease of that difference.
The spatial error common to the two estimates
cancels in the comparison, leaving the finite-depth
bias and time-rounding effects.
It is the errors against the independent continuum
reference that establish spatial convergence.
The tail calculation for this particular production
profile uses its known finite Fourier support;
the metric anchor remains part of the supplied data.


\section{Consequences for finite-window reconstruction}
The information in two Gowdy slices depends on the measured
components, background, and placement of the observations.
The homogeneous fold and exact rank-two collisions give
specific failures. Null transport, and in some cases the
metric momentum constraint, supply information at other
placements. The nonlinear estimates quantify when the
missing velocity can be stably recovered without supplying
it directly. The finite-grid certificate explains the
archived reconstruction within its precise numerical model.
The metric-modulus estimate additionally identifies which
regularity is needed to recover an absolute limiting
quantity once a metric anchor is supplied.

\appendix
\section{An extension to nontrivial torus bundles}
\label{sec:twisted}
The Gowdy-to-Ernst examples on the three-torus with
matched full metrics use a periodic transformed $Q$,
so their scalar seed has zero-mean time velocity.
On a nontrivial torus bundle a nonzero current period is
admissible. This yields a rank-two collision for the
original four-field record with asymptotic speed
strictly between zero and one at every angle.

For an integer $n\ne0$ impose the type-II identification
\[
 (\theta,x,y)\sim(\theta+2\pi,x-ny,y),
 \qquad x,y\in\mathbb R/\mathbb Z.
\]
Take periodic $P,\lambda$ and
$Q(t,\theta+2\pi)=Q(t,\theta)+n$.
The orbit form $dx+Q\,dy$ descends under the displayed
identification. This is the standard unipotent
monodromy in twisted Gowdy \cite{rendall-twisted}.

\begin{proposition}
\label{prop:twisted}
For $n=1$ there are two distinct locally Gowdy vacuum
metrics with this compact torus-bundle topology, having
rank-two target maps on open sets and the same
$P(a),Q(a),Q_t(a),P(b)$ and initial metric anchor
in the common marked bundle coordinates.
Their asymptotic speeds are uniformly in $(0,1)$.
In the normalized current-one family generated from a
constant initial seed position, the terminal metric
potential distinguishes the two branches if the common
angular current period is nonzero.
\end{proposition}

\begin{proof}
Take the two geodesics from the subcritical row
of Table~\ref{tab:collisions}. Their speeds differ,
but they have common $q>0$ and common $P$ value
at both phase values $0$ and $\ell=3$.
Translate the initial $P$ to $P_0$ and rescale $Q$
by $e^{-P_0}$ to preserve the orthonormal components
of their initial velocities. Choose
$C=q e^{P_0}=1/(2\pi)$.
Use the phase
\[
 w(t,\theta)=\log(t/a)+\epsilon u_{k_0}(t)\cos(k_0\theta),
 \qquad u_{k_0}(a)=u_{k_0}(b)=0,\quad b/a=e^3 .
\]
Both seed metric potentials are periodic because
the initial orbit positions are constant, so the
periodic momentum constraint holds.
Apply \eqref{eq:ernst}. The transformed angular
period of $Q$ is
$C\int_0^{2\pi}t w_t\,d\theta=1$,
the required monodromy. The transformed metric potential
is periodic and the transformed $Q$ can be chosen common.
At both nodes $w_\theta=0$, and the initial
$Q_t$ and the measured $P$ values therefore agree.
The initial metric anchor can be taken common.
The missing initial velocities differ by a nonzero
multiple of $1+\epsilon\cos(k_0\theta)$.
The rank formula of Theorem~\ref{th:rank-two} still
holds at $a$, since $w_\theta(a)=0$.
Its factor $1+aP_w w_t$ is nonzero for both branches
for sufficiently small amplitudes: at zero amplitude
it is $1+g_j>0$.

The limit $t w_t\to1+\epsilon J_0(k_0a)\cos(k_0\theta)$
is positive everywhere for small $\epsilon$.
Then $w\to-\infty$ uniformly and $P_w\to-\sigma_j$
on each seed orbit. The transverse contribution
$t e^{\widetilde P}\widetilde Q_t$
vanishes as $t\downarrow0$.
The transformed speed is therefore
$1-\sigma_j(1+\epsilon J_0(k_0a)\cos(k_0\theta))$,
strictly within $(0,1)$ for both subcritical speeds
in the selected row and small enough amplitude.

For the final claim, normalize the seed current as
in Proposition~\ref{prop:ernst-fiber}.
The initial transformed $Q$ fixes the entire scalar
phase $\psi$, including the mean drift.
The energy comparison at the common angular origin
uses the integral $I$ in that proposition, positive
there whenever the common terminal phase is
nonzero. If terminal $P$ agrees, equality of terminal
$\lambda$ fixes the common $\kappa^2$ because the
$-2P-\log t$ terms in the transform agree.
Equality of terminal $P$ and \eqref{eq:ernst-fiber}
then determine $\zeta$ wherever $\psi(b)\ne0$.
For nonzero monodromy the mean terminal phase is
nonzero as $b>a$, so such an angle exists.
\end{proof}


\clearpage
\begin{thebibliography}{99}
\setlength{\itemsep}{1pt}
\bibitem{aldroubi} A. Aldroubi, C. Cabrelli, U. Molter and S. Tang,
Dynamical sampling, \emph{Appl. Comput. Harmon. Anal.} 42 (2017),
378--401. \href{https://arxiv.org/abs/1409.8333}{arXiv:1409.8333}.

\bibitem{amorim} P. Amorim, C. Bernardi and P. G. LeFloch,
Computing Gowdy spacetimes via spectral evolution in future and past
directions, \emph{Class. Quantum Grav.} 26 (2009), 025007.
\href{https://arxiv.org/abs/0811.1266}{arXiv:0811.1266}.

\bibitem{berger} B. K. Berger and V. Moncrief,
Numerical investigation of cosmological singularities,
\emph{Phys. Rev. D} 48 (1993), 4676--4687.
\href{https://arxiv.org/abs/gr-qc/9307032}{arXiv:gr-qc/9307032},
\S\S III--IV, especially Eq.~(3.19).

\bibitem{cssw} J. Christensen, F. Gonzalez, T. Kakehi and J. Wang,
The snapshot problem for the wave equation,
\emph{Adv. Math.} 444 (2024), 109629.
\href{https://arxiv.org/abs/2308.12208}{arXiv:2308.12208}.

\bibitem{dlr} B. Dehman, J. Le Rousseau and M. Léautaud,
Controllability of two coupled wave equations on a compact manifold,
\emph{Arch. Ration. Mech. Anal.} 211 (2014), 113--187.
\href{https://doi.org/10.1007/s00205-013-0670-4}{DOI}.

\bibitem{clw} Y. Cui, C. Laurent and Z. Wang,
On the observability inequality of coupled wave equations:
the case without boundary,
\emph{ESAIM Control Optim. Calc. Var.} 26 (2020), 14.
\href{https://doi.org/10.1051/cocv/2020004}{DOI}.

\bibitem{cindea} N. C\^{\i}ndea, A. Imperiale and P. Moireau,
Data assimilation of time under-sampled measurements using
observers, the wave-like equation example,
\emph{ESAIM Control Optim. Calc. Var.} 21 (2015), 635--669.
\href{https://doi.org/10.1051/cocv/2014042}{DOI}.

\bibitem{dlmf} NIST Digital Library of Mathematical Functions,
\href{https://dlmf.nist.gov/10.17}{\S 10.17, Asymptotic Expansions for Large Argument}.

\bibitem{grubisic} B. Grubi\v{s}i\'{c} and V. Moncrief,
Asymptotic behavior of the $T^3\times\mathbb R$ Gowdy space-times,
\emph{Phys. Rev. D} 47 (1993), 2371--2382.
\href{https://arxiv.org/abs/gr-qc/9209006}{arXiv:gr-qc/9209006}.

\bibitem{li} W. Li,
Scattering towards the singularity for the wave equation and
the linearized Einstein--scalar field system in Kasner spacetimes.
\href{https://arxiv.org/abs/2401.08437}{arXiv:2401.08437}.

\bibitem{libounces} W. Li,
BKL bounces outside homogeneity: Gowdy symmetric spacetimes.
\href{https://arxiv.org/abs/2408.12427}{arXiv:2408.12427}.

\bibitem{lim} W. C. Lim, New explicit spike solutions---non-local component
of the generalized Mixmaster attractor,
\emph{Class. Quantum Grav.} 25 (2008), 045014.
\href{https://arxiv.org/abs/0710.0628}{arXiv:0710.0628}, Eq.~(22).

\bibitem{old} R. Bibb, Boundary-State Sufficiency in Gowdy Cosmology.
\href{https://doi.org/10.5281/zenodo.21959052}{Zenodo preprint (2026)}.

\bibitem{rendall-twisted} A. D. Rendall,
Dynamics of solutions of the Einstein equations with twisted
Gowdy symmetry, \emph{J. Geom. Phys.} 62 (2012), 569--577.
\href{https://doi.org/10.1016/j.geomphys.2011.04.016}{DOI}.

\bibitem{ringstrom} H. Ringstr\"om, On a wave map equation arising in
general relativity, \emph{Comm. Pure Appl. Math.} 57 (2004), 657--703.
\href{https://doi.org/10.1002/cpa.20015}{DOI}.

\bibitem{ringstromavtd} H. Ringstr\"om, Asymptotic expansions close
to the singularity in Gowdy spacetimes,
\emph{Class. Quantum Grav.} 21 (2004), S305--S322.
\href{https://doi.org/10.1088/0264-9381/21/3/019}{DOI}.

\bibitem{rw} A. D. Rendall and M. Weaver,
Manufacture of Gowdy spacetimes with spikes,
\emph{Class. Quantum Grav.} 18 (2001), 2959--2976.
\href{https://arxiv.org/abs/gr-qc/0103102}{arXiv:gr-qc/0103102}.

\bibitem{vest} A. Vest, Observation of vibrating systems at different time
instants, \emph{Electron. J. Qual. Theory Differ. Equ.} 2014, no.~14, 1--15.
\href{https://arxiv.org/abs/1305.6461}{arXiv:1305.6461}.
\end{thebibliography}
\end{document}
